\section{AMC 12 12A 2022}
        \begin{enumerate}
        	\item (\textit{AMC 12 12A 2022}) What is the value of 
 \[3+\frac{1}{3+\frac{1}{3+\frac13}}?\]
$\textbf{(A)}\ \frac{31}{10}\qquad\textbf{(B)}\ \frac{49}{15}\qquad\textbf{(C)}\ \frac{33}{10}\qquad\textbf{(D)}\ \frac{109}{33}\qquad\textbf{(E)}\ \frac{15}{4}$

 \textbf{Answer: }D



\textbf{Solution 1}

We have 
 \begin{align*} 3+\frac{1}{3+\frac{1}{3+\frac13}} &= 3+\frac{1}{3+\frac{1}{\left(\frac{10}{3}\right)}} \\ &= 3+\frac{1}{3+\frac{3}{10}} \\ &= 3+\frac{1}{\left(\frac{33}{10}\right)} \\ &= 3+\frac{10}{33} \\ &= \boxed{\textbf{(D)}\ \frac{109}{33}}. \end{align*}
~MRENTHUSIASM

 



\textbf{Solution 2}

Continued fractions with integer parts $q_i$ and numerators all $1$ can be calculated as 
 \begin{align*} \dfrac{[q_0,q_1,q_2,\ldots,q_n]}{[q_1,q_2,\ldots,q_n]} \end{align*}
where
 \begin{align*} []&=1 \\ [q_0,q_1,q_2,\ldots,q_n]&=q_0[q_1,q_2,\ldots,q_n]+[q_2,\ldots,q_n]\\ \end{align*}





\textbf{Solution 3}

Finite continued fractions of form $n+\frac{1}{n+\frac{1}{n+\cdots}}=\frac{x}{y}$ have linear combinations of $x, y$ that solve Pell's Equation. Specifically, the denominator $y$ and numerator $x$ are solutions to the Diophantine equation $(n^2+4)\left(\frac{y}{2}\right)^2-\left(x-\frac{ny}{2}\right)^2=\pm{1}$. So for this problem in particular, the denominator $y$ and numerator $x$ are solutions to the Diophantine equation $13\left(\frac{y}{2}\right)^2-\left(x-\frac{3y}{2}\right)^2=\pm{1}$. That leaves two answers. Since the number of $1$'s in the continued fraction is odd, we further narrow it down to $13\left(\frac{y}{2}\right)^2-\left(x-\frac{3y}{2}\right)^2=-1$, which only leaves us with $1$ answer and that is $(x,y)=(109,33)$ which means $\boxed{\textbf{(D)}\ \frac{109}{33}}$. 

 (Note: Integer solutions increase exponentially, so our next solution will have a numerator greater than $3^2(109)$. Therefore, when you don't see numerators greater than $3^2(109)$ in the answer choices, this method should be fine.)

 


	\item (\textit{AMC 12 12A 2022}) The sum of three numbers is $96.$ The first number is $6$ times the third number, and the third number is $40$ less than the second number. What is the absolute value of the difference between the first and second numbers?


 $\textbf{(A) } 1 \qquad \textbf{(B) } 2 \qquad \textbf{(C) } 3 \qquad \textbf{(D) } 4 \qquad \textbf{(E) } 5$

 \textbf{Answer: }E



\textbf{Solution 1}

Let $x$ be the third number. It follows that the first number is $6x,$ and the second number is $x+40.$

 We have 
 \[6x+(x+40)+x=8x+40=96,\] from which $x=7.$

 Therefore, the first number is $42,$ and the second number is $47.$ Their absolute value of the difference is $|42-47|=\boxed{\textbf{(E) } 5}.$




\textbf{Solution 2}

Solve this using a system of equations. Let $x,y,$ and $z$ be the three numbers, respectively. We get three equations:
 \begin{align*} x+y+z&=96, \\ x&=6z, \\ z&=y-40. \end{align*}
Rewriting the third equation gives us $y=z+40,$ so we can substitute $x$ as $6z$ and $y$ as $z+40.$

 Therefore, we get
 \begin{align*} 6z+(z+40)+z&=96 \\ 8z+40&=96 \\ 8z&=56 \\ z&=7. \end{align*}
Substituting 7 in for $z$ gives us $x=6z=6(7)=42$ and $y=z+40=7+40=47.$

 So, the answer is $|x-y|=|42-47|=\boxed{\textbf{(E) } 5}.$




\textbf{Solution 3}

In accordance with Solution 2, 
 \[y = z+40, x = 6z \implies |x-y| = |6z - z - 40| = 5|z - 8| \implies \boxed{\textbf{(E) } 5}.\]\textbf{vladimir.shelomovskii@gmail.com, vvsss}

 


	\item (\textit{AMC 12 12A 2022}) Five rectangles, $A$, $B$, $C$, $D$, and $E$, are arranged in a square as shown below. These rectangles have dimensions $1\times6$, $2\times4$, $5\times6$, $2\times7$, and $2\times3$, respectively. (The figure is not drawn to scale.) Which of the five rectangles is the shaded one in the middle?

                    \begin{figure}[H]
                        \centering
                        \includegraphics[width=0.5\textwidth]{Problem/AMC_12_12A_2022_3_Problem_Figure_1.png}
                    \end{figure}
                    
$\textbf{(A) }A\qquad\textbf{(B) }B \qquad\textbf{(C) }C \qquad\textbf{(D) }D\qquad\textbf{(E) }E$

 \textbf{Answer: }B



\textbf{Solution 1 (Area and Perimeter of Square)}

The area of this square is equal to $6 + 8 + 30 + 14 + 6 = 64$, and thus its side lengths are $8$. The sum of the dimensions of the rectangles are $2 + 7 + 5 + 6 + 2 + 3 + 1 + 6 + 2 + 4 = 38$. Thus, because the perimeter of the square is $32$, the rectangle on the inside must have a perimeter of $6 \cdot 2 = 12$. The only rectangle that works is $\boxed{\textbf{(B) }B}$.

 



\textbf{Solution 2 (Perimeter of Square)}

Note that the perimeter of the square is the sum of four pairs of dimensions (eight values are added). Adding up all dimensions gives us $2+7+5+6+2+3+1+6+2+4=38$. We know that as the square's side length is an integer, the perimeter must be divisible by $4$. Testing out by subtracting all five pairs of dimensions from $38$, only $2\times4$ works since $38-2-4=32=8\cdot4$, which corresponds with $\boxed{\textbf{(B) }B}$.

 



\textbf{Solution 3 (Observations)}

Note that rectangle $D$ must be on the edge. Without loss of generality, let the top-left rectangle be $D,$ as shown below:
                    \begin{figure}[H]
                        \centering
                        \includegraphics[width=0.5\textwidth]{Solution/AMC_12_12A_2022_3_Solution_Figure_1.png}
                    \end{figure}
                    It is clear that $x=1,$ so we can determine Rectangle $A.$

 Continuing with a similar process, we can determine Rectangles $C,E,$ and $B,$ in this order. The answer is $\boxed{\textbf{(B) }B}$ as shown below.
                    \begin{figure}[H]
                        \centering
                        \includegraphics[width=0.5\textwidth]{Solution/AMC_12_12A_2022_3_Solution_Figure_2.png}
                    \end{figure}
                    ~MRENTHUSIASM

 



\textbf{Solution 4 (Observations)}

Let's label some points:
                    \begin{figure}[H]
                        \centering
                        \includegraphics[width=0.5\textwidth]{Solution/AMC_12_12A_2022_3_Solution_Figure_3.png}
                    \end{figure}
                    By finding the dimensions of the middle rectangle, we need to find the dimensions of the other four rectangles. We have a rule: 
 \[AB + BC = CD + DE = EF + FG = GH + AH.\]Let's make a list of all dimensions of the rectangles from the diagram. We have to fill in the dimensions from up above, but still apply to the rule:
 \begin{align*} AB&\times AH \\ CD&\times BC \\ EF&\times DE \\ GH&\times FG \end{align*}
By applying the rule, we get $AB=2, BC=6, CD=5, DE=3, EF=2, FG=6, GH=1$, and $AH=7$.

 By substitution, we get this list
 \begin{align*} 2&\times 7 \\ 5&\times 6 \\ 2&\times 3 \\ 1&\times 6 \\ \end{align*}
(This also tells us that the diagram is not drawn to scale.)

 Notice how the only dimension not used in the list was $2\times 4$ and that corresponds with B so the answer is, $\boxed{\textbf{(B) }B}.$



	\item (\textit{AMC 12 12A 2022}) The least common multiple of a positive integer $n$ and $18$ is $180$, and the greatest common divisor of $n$ and $45$ is $15$. What is the sum of the digits of $n$?


 $\textbf{(A) } 3 \qquad \textbf{(B) } 6 \qquad \textbf{(C) } 8 \qquad \textbf{(D) } 9 \qquad \textbf{(E) } 12$

 \textbf{Answer: }B



\textbf{Solution 1}

Note that
 \begin{align*} 18 &= 2\cdot3^2, \\ 180 &= 2^2\cdot3^2\cdot5, \\ 45 &= 3^2\cdot5 \\ 15 &= 3\cdot5. \end{align*}
Let $n = 2^a\cdot3^b\cdot5^c.$ It follows that:

 
            \begin{enumerate} 
        \item From the least common multiple condition, we have 
 \[\operatorname{lcm}(n,18) = \operatorname{lcm}(2^a\cdot3^b\cdot5^c,2\cdot3^2) = 2^{\max(a,1)}\cdot3^{\max(b,2)}\cdot5^{\max(c,0)} = 2^2\cdot3^2\cdot5,\] from which $a=2, b\in\{0,1,2\},$ and $c=1.$

        \item From the greatest common divisor condition, we have 
 \[\gcd(n,45) = \gcd(2^2\cdot3^b\cdot5,3^2\cdot5) = 2^{\min(2,0)}\cdot3^{\min(b,2)}\cdot5^{\min(1,1)} = 3\cdot5,\] from which $b=1.$

            \end{enumerate}
            Together, we conclude that $n=2^2\cdot3\cdot5=60.$ The sum of its digits is $6+0=\boxed{\textbf{(B) } 6}.$




\textbf{Solution 2}

The options for $\text{lcm}(x, 18)=180$ are $20$, $60$, and $180$. The options for $\text{gcd}(y, 45)=15$ are $15$, $30$, $60$, $75$, etc. We see that $60$ appears in both lists; therefore, $6+0=\boxed{\textbf{(B) } 6}$.

 \textbf{Remark} If you ignore or mess up the LCM, and get  $n=15$, you'll still get the correct answer.

 


	\item (\textit{AMC 12 12A 2022}) The <em>taxicab distance</em> between points $(x_1, y_1)$ and $(x_2, y_2)$ in the coordinate plane is given by 
 \[|x_1 - x_2| + |y_1 - y_2|.\]
For how many points $P$ with integer coordinates is the taxicab distance between $P$ and the origin less than or equal to $20$?


 $\textbf{(A)} \, 441 \qquad\textbf{(B)} \, 761 \qquad\textbf{(C)} \, 841 \qquad\textbf{(D)} \, 921  \qquad\textbf{(E)} \, 924$

 \textbf{Answer: }C



\textbf{Solution 1 (Arithmetic Series)}

Let us consider the number of points for a certain $x$-coordinate. For any $x$, the viable points are in the range $[-20 + |x|, 20 - |x|]$. This means that our total sum is equal to 
 \begin{align*} 1 + 3 + 5 + \cdots + 41 + 39 + 37 + \cdots + 1 &= (1 + 3 + 5 + \cdots + 39) + (1 + 3 + 5 + \cdots + 41) \\ & = 20^2 + 21^2 \\ & = 29^2 \\ &= \boxed{\textbf{(C)} \, 841}. \end{align*}
~mathboy100

 



\textbf{Solution 2 (Arithmetic Series With Symmetry)}

Since the second point is the origin, this is equivalent to finding all points $(x, y)$ such that $|x| + |y| \leq 20$. Due to the absolute values, the set of all such points will be symmetric about the origin meaning we can focus on the first quadrant and multiply by $4$.

 To avoid overcounts, ignore points on the axes. This means $x, y > 0$. If $x = 1$, there are $19$ solutions for $y$ ($y = 1, 2, 3, \ldots, 19$). If $x = 2$, there are $18$ solutions. This pattern repeats until $x = 19$, at which point there is $1$ solution for $y$.

 So we get $19 + 18 + 17 + \cdots + 1 = \frac{19(20)}{2} = 190$ points in the first quadrant. Multiplying by $4$ gives $760$. Now, the $x$ axis has $y = 0$ which gives $|x| \leq 20$, meaning there are $41$ solutions. This is the same with the $y$ axis, but we overcounted the origin by $1$.

 Our final answer is $760 + 41 + 41 - 1 = \boxed{\textbf{(C)} \, 841}$.

 



\textbf{Solution 3 (Triangular Numbers)}

This solution refers to the \textbf{Diagram} section.
                    \begin{figure}[H]
                        \centering
                        \includegraphics[width=0.5\textwidth]{Solution/AMC_12_12A_2022_5_Solution_Figure_1.png}
                    \end{figure}
                    The problem can be visualized as depicted on the right split equally into four "triangular" parts excluding the origin. The "triangular" parts are identical the ones that would be used in a visual proof of the formula for triangular numbers. Becuase of this the number of points in each part is equal to $\frac{n(n+1)}{2}$ where $n$ is the length of a "leg" of the "triangle" which is $20$ for this problem. Substituting and computing, we get $210.$ Multiplying by $4$ and adding $1$ to account for all parts and the origin, we get $210\cdot4 + 1 = \boxed{\textbf{(C)} \, 841}.$




\textbf{Solution 4 (Two Square Arrays)}

This solution refers to the \textbf{Diagram} section.

 As shown below, the taxicab distance between each red point and the origin is even, and the taxicab distance between each blue point and the origin is odd.
                    \begin{figure}[H]
                        \centering
                        \includegraphics[width=0.5\textwidth]{Solution/AMC_12_12A_2022_5_Solution_Figure_2.png}
                    \end{figure}
                    Note that the red array consists of $21^2=441$ points, and the blue array consists of $20^2=400$ points.

 Together, the answer is $441+400=\boxed{\textbf{(C)} \, 841}.$




\textbf{Solution 5 (Pick's Theorem)}

Let $P = (x, y)$. Since the problem asks for taxicab distances from the origin, we want $|x| + |y| \le 20$. The graph of all solutions to this equation on the $xy$-plane is a square with vertices at $(0, \pm 20)$ and $(\pm 20, 0)$ (In order to prove this, one can divide the sections of this graph into casework on the four quadrants, and tie together the resulting branches.) We want the number of lattice points on the border of the square and inside the square. Each side of the square goes through an equal number of lattice points, so if we focus on one side going from $(0,20)$ to $(20, 0)$, we can see that it goes through $21$ points in total. In addition, each of the vertices gets counted twice, so the total number of border points is $21\cdot4 - 4 = 80$. Also, the area of the square is $800$, so when we plug this information inside Pick's theorem, we get $800 = i + \frac{80}{2} - 1 \implies i = 761$. Then our answer is $761+80 = \boxed{\textbf{(C)} \, 841}.$




\textbf{Solution 6 (Stars and Bars)}

Instead of considering all points with integer coordinates, first consider points with nonnegative coordinates only. Then, we want $x + y \le 20$ where $x$ and $y$ are nonnegative integers. We can introduce a third variable, say $z$, such that $z = 20 - (x + y)$. Note that counting the number of ways to have $x + y + z = 20$ is the same as counting the number of ways to have $x + y \le 20$. Therefore, by stars and bars, there are $\dbinom{20 + 3 - 1}{3 - 1} = 231$ solutions with nonnegative integer coordinates.

 Then, we can copy our solutions over to the other four quadrants. First, so as not to overcount, we remove all points on the axes. There are $20 + 20 + 1 = 41$ such points with nonnegative integer coordinates. We multiply the $190$ remaining points by $4$ to get $760$ points that are not on the axes. Then, we can add back the $41$ nonnegative points on the axes, as well as the $40$ other points on the negative axes to get $760 + 41 + 40 = \boxed{\textbf{(C)} \, 841}.$




\textbf{Solution 7 (Fast, Cheese)}

The number of lattice points is roughly equal to the area of a rhombus with diagonals of length 41. The area of the quadrilateral is $41^2/2 = 840.5$. The closest answer is $\boxed{\textbf{(C)} \, 841}.$.

 


	\item (\textit{AMC 12 12A 2022}) A data set consists of $6$ (not distinct) positive integers: $1$, $7$, $5$, $2$, $5$, and $X$. The
average (arithmetic mean) of the $6$ numbers equals a value in the data set. What is
the sum of all positive values of $X$?


 $\textbf{(A) } 10 \qquad \textbf{(B) } 26 \qquad \textbf{(C) } 32 \qquad \textbf{(D) } 36 \qquad \textbf{(E) } 40$

 \textbf{Answer: }D



\textbf{Solution (Casework)}

First, note that $1+7+5+2+5=20$. There are $3$ possible cases:

 \textbf{Case 1: the mean is $5$.}

 $X = 5 \cdot 6 - 20 = 10$.

 \textbf{Case 2: the mean is $7$.}

 $X = 7 \cdot 6 - 20 = 22$.

 \textbf{Case 3: the mean is $X$.}

 $X= \frac{20+X}{6} \Rightarrow X=4$.

 Therefore, the answer is $10+22+4=\boxed{\textbf{(D) }36}$.

 


	\item (\textit{AMC 12 12A 2022}) A rectangle is partitioned into $5$ regions as shown. Each region is to be painted a solid color - red, orange, yellow, blue, or green - so that regions that touch are painted different colors, and colors can be used more than once. How many different colorings are possible?


 
                    \begin{figure}[H]
                        \centering
                        \includegraphics[width=0.5\textwidth]{Problem/AMC_12_12A_2022_7_Problem_Figure_1.png}
                    \end{figure}
                    
$\textbf{(A) }120\qquad\textbf{(B) }270\qquad\textbf{(C) }360\qquad\textbf{(D) }540\qquad\textbf{(E) }720$

 \textbf{Answer: }D



\textbf{Solution 1}

The top left rectangle can be $5$ possible colors. Then the bottom left region can only be $4$ possible colors, and the bottom middle can only be $3$ colors since it is next to the top left and bottom left. Similarly, we have $3$ choices for the top right and $3$ choices for the bottom right, which gives us a total of $5\cdot4\cdot3\cdot3\cdot3=\boxed{\textbf{(D) }540}$.

 Note that the number of choices for a rectangle is $5-n$ when there are $n$ rectangles adjacent to it. ~GoldemGem4848

 



\textbf{Solution 2 (Casework)}

Case 1: All the rectangles are different colors. It would be $5! = 120$ choices. 

 Case 2: Two rectangles that are the same color. Grouping these two rectangles as one gives us $5\cdot4\cdot3\cdot2 = 120$. But, you need to multiply this number by three because the same-colored rectangles can be chosen at the top left and bottom right, the top right and bottom left, or the bottom right and bottom left, which gives us a grand total of $360$.

 Case 3: We have two sets of rectangles chosen from these choices (top right & bottom left, top left & bottom right) that have the same color. However, the choice of the bottom left and bottom right does not work for this case, as the second pair would be chosen from two touching rectangles. Again, grouping the same-colored rectangles gives us $5\cdot4\cdot3 = 60$. 

 Therefore, we have $120 + 360 + 60 = \boxed{\textbf{(D) }540}$.

 



\textbf{Solution 3 (Different Kind Of Casework)}

Let us first choose the color of the bottom-middle rectangle, as it is touching all the other rectangles and has to be a different color than all of them. This rectangle has $5$ possibilities. To make things easier for us, let's assume we chose the color red. Now, lets choose the color of the top-left rectangle. This has $4$ options: orange, yellow, blue, or green (it can't be the same color as the middle rectangle, which we assumed was red). Lets assume this is orange. Now, we know that the bottom-left rectangle only has $3$ choices: yellow, blue, or green. Lets assume this is colored yellow. Now, lets do casework on the color of the bottom right-rectangle. It either can be the same color as the top-left (orange), or different.CASE 1: (Its the same color)We then know that the bottom-right rectangle is also colored orange. We now know that the top-left rectangle only has $3$ choices: yellow, blue, or green (it can't be orange or red). The sum of the possibilities of this case is $5\cdot4\cdot3\cdot3$.

 CASE 2: (Its a different color)We then know that the bottom-right rectangle has $3$ options (yellow, blue, or green, since it can't be red, and it can't be orange since we already covered that in the case 1). Lets assume it is colored yellow. We know have $3$ choices for the top-right rectangle (blue or green). The sum of the possibilities of this case is $5\cdot4\cdot3\cdot3\cdot2$

 We can then sum up the possibilities of our two cases, which is $5\cdot4\cdot3\cdot3+5\cdot4\cdot3\cdot3\cdot2 = 5\cdot4\cdot3(3+3\cdot2) = 60\cdot9 = \boxed{\textbf{(D) }540}$



	\item (\textit{AMC 12 12A 2022}) The infinite product

 \[\sqrt[3]{10} \cdot \sqrt[3]{\sqrt[3]{10}} \cdot \sqrt[3]{\sqrt[3]{\sqrt[3]{10}}} \cdots\]
evaluates to a real number. What is that number?


 $\textbf{(A) }\sqrt{10}\qquad\textbf{(B) }\sqrt[3]{100}\qquad\textbf{(C) }\sqrt[4]{1000}\qquad\textbf{(D) }10\qquad\textbf{(E) }10\sqrt[3]{10}$

 \textbf{Answer: }A



\textbf{Solution 1}

We can write $\sqrt[3]{10}$ as $10 ^ \frac{1}{3}$. Similarly, $\sqrt[3]{\sqrt[3]{10}} = (10 ^ \frac{1}{3}) ^ \frac{1}{3} = 10 ^ \frac{1}{3^2}$.

 By continuing this, we get the form
 \[10 ^ \frac{1}{3} \cdot 10 ^ \frac{1}{3^2} \cdot 10 ^ \frac{1}{3^3} \cdots,\]which is
 \[10 ^ {\frac{1}{3} + \frac{1}{3^2} + \frac{1}{3^3} + \cdots}.\]Using the formula for an infinite geometric series $S = \frac{a}{1-r}$, we get
 \[\frac{1}{3} + \frac{1}{3^2} + \frac{1}{3^3} + \cdots = \frac{\frac{1}{3}}{1-\frac{1}{3}} = \frac{1}{2}.\]Thus, our answer is $10 ^ \frac{1}{2} = \boxed{\textbf{(A) }\sqrt{10}}$.

 - phuang1024

 



\textbf{Solution 2}

We can write this infinite product as $L$ (we know from the answer choices that the product must converge):
 \[L = \sqrt[3]{10} \cdot \sqrt[3]{\sqrt[3]{10}} \cdot \sqrt[3]{\sqrt[3]{\sqrt[3]{10}}} \cdots.\]If we raise everything to the third power, we get:
 \[L^3 =  10 \, \cdot \, \sqrt[3]{10} \, \cdot \, \sqrt[3]{\sqrt[3]{10}} \cdots = 10L \implies L^3 - 10L = 0 \implies L \in \left\{0, \pm \sqrt{10}\right\}.\]Since $L$ is positive (as it is an infinite product of positive numbers), it must be that $L = \boxed{\textbf{(A) }\sqrt{10}}.$




\textbf{Solution 3}

Move the first term inside the second radical. We get
 \[\sqrt[3]{10} \cdot \sqrt[3]{\sqrt[3]{10}} \cdot \sqrt[3]{\sqrt[3]{\sqrt[3]{10}}} \cdots = \sqrt[3]{10\sqrt[3]{10}} \cdot \sqrt[3]{\sqrt[3]{\sqrt[3]{10}}} \cdots.\]Do this for the third radical as well:
 \[\sqrt[3]{10\sqrt[3]{10}} \cdot \sqrt[3]{\sqrt[3]{\sqrt[3]{10}}} \cdots = \sqrt[3]{10\sqrt[3]{10}\sqrt[3]{\sqrt[3]{10}}} \cdots = \sqrt[3]{10\sqrt[3]{10\sqrt[3]{10\cdots}}}.\]It is clear what the pattern is. Setting the answer as $P,$ we have 
 \[P = \sqrt[3]{10P},\] from which $P = \boxed{\sqrt{10}}.$




\textbf{Solution 4}

Set the product equal to P. We get
 \[\sqrt[3]{10} \cdot \sqrt[3]{\sqrt[3]{10}} \cdot \sqrt[3]{\sqrt[3]{\sqrt[3]{10}}} \cdots=P\] Since this is an infinite product, there may exist a clever manipulation where we set two different espressions involving $P$ equal, from which we could solve for a number. Since all terms are raised to the $\frac{1}{3}$ power, we can cube both sides of our equation. We would get 
 \[10 \cdot \sqrt[3]{10} \cdot \sqrt[3]{\sqrt[3]{10}} \cdot \sqrt[3]{\sqrt[3]{\sqrt[3]{10}}} \cdots=P^3\] From this, we know that $\frac{P^3}{10}=P$, and to this equation, the only solution is $P = \boxed{\sqrt{10}}$



	\item (\textit{AMC 12 12A 2022}) On Halloween $31$ children walked into the principal's office asking for candy. They
can be classified into three types: Some always lie; some always tell the truth; and
some alternately lie and tell the truth. The alternaters arbitrarily choose their first
response, either a lie or the truth, but each subsequent statement has the opposite
truth value from its predecessor. The principal asked everyone the same three
questions in this order.


 "Are you a truth-teller?" The principal gave a piece of candy to each of the $22$
children who answered yes.


 "Are you an alternater?" The principal gave a piece of candy to each of the $15$
children who answered yes.


 "Are you a liar?" The principal gave a piece of candy to each of the $9$ children who
answered yes.


 How many pieces of candy in all did the principal give to the children who always
tell the truth?


 $\textbf{(A) } 7 \qquad \textbf{(B) } 12 \qquad \textbf{(C) } 21 \qquad \textbf{(D) } 27 \qquad \textbf{(E) } 31$

 \textbf{Answer: }A



\textbf{Solution 1}

Note that:

 
            \begin{itemize}
        \item Truth-tellers would answer yes-no-no to the three questions in this order.
            \end{itemize}
            
            \begin{itemize}
        \item Liars would answer yes-yes-no to the three questions in this order.
            \end{itemize}
            
            \begin{itemize}
        \item Alternaters who responded truth-lie-truth would answer no-no-no to the three questions in this order.
            \end{itemize}
            
            \begin{itemize}
        \item Alternaters who responded lie-truth-lie would answer yes-yes-yes to the three questions in this order.
            \end{itemize}
            Suppose that there are $T$ truth-tellers, $L$ liars, and $A$ alternaters who responded lie-truth-lie.

 The conditions of the first two questions imply that
 \begin{align*} T+L+A&=22, \\ L+A&=15. \end{align*}
Subtracting the second equation from the first, we have $T=22-15=\boxed{\textbf{(A) } 7}.$
<u>\textbf{Remark}</u>

 The condition "$31$ children" and the third question are extraneous. However, they imply $A=9$ and $L=6,$ from which there are $31-T-L-A=9$ alternaters who responded truth-lie-truth.

 



\textbf{Solution 2}

Consider when the principal asks "Are you a liar?": The truth tellers truthfully say no, and the liars lie and say no. This leaves only alternaters who lie on this question to answer yes. Thus, all $9$ children that answered yes are alternaters that falsely answer Questions 1 and 3, and truthfully answer Question 2. The rest of the alternaters, however many there are, have the opposite behavior.

 Consider the second question, "Are you an alternater?": The truth tellers again answer no, the liars falsely answer yes, and alternaters that truthfully answer also say yes. From the previous part, we know that $9$ alternaters truthfully answer here. Because only liars and $9$ alternaters answer yes, we can deduce that there are $15-9=6$ liars.

 Consider the first question, "Are you a truth teller?": Truth tellers say yes, liars also say yes, and alternaters that lie on this question also say yes. From the first part, we know that $9$ alternaters lie here. From the previous part, we know that there are $6$ liars. Because only the number of truth tellers is unknown here, we can deduce that there are $22-9-6=7$ truth tellers.

 The final question is how many pieces of candy did the principal give to truth tellers. Because truth tellers answer yes on only the first question, we know that all $7$ of them said yes once, resulting in $\boxed{\textbf{(A) } 7}$ pieces of candy.

 


 



\textbf{Solution 3 (case analysis)}

There are $T$ truthtellers, $L$ liars, and $A$ alternators. We translate from logic to math via \textbf{formalization}.

 From statement 3 we see that $A \ge 9$. From statement 2 we have that $A+L = 24$, and from statement 1 we have that $T+L = 13$.

 We also know that $T+L+A = 31$. Having a system of three, we see that case 3 is redundant and unnecessary. Therefore, solving
 \[T+L+A=31\]
 \[T+L=13\]
 \[A+L=24\]

 the solution $(T,A,L) = (\boxed{7},18,6)$. This is exactly solution 1 and 2 except faster and mathematically intuitive rather than logically intuitive. 

 



\textbf{Solution 4 (logical deduction via answer choices)}

Note that liars will always lie and truthtellers will always tell the truth. Therefore, when he asks, "Are you a liar", only the alternators who answered "Yes" to question 1 will answer question 3 as "Yes". Therefore, we currently know 9 alternators exist.

 Now when he asks, "Are you are an alternator", we know the remaining alternators and the liars will both say "Yes". Therefore, we know that the remaining alternators and liars answer yes to that. In addition, when he asks "Are you a truthteller", all the 9 alternators will say yes, so we have 13 people, some of which are truthtellers, and some of which are liars. This reduces our answer to either 7 or 12.

 Now, say we have 7 truthtellers. Then, that must mean we have 6 liars. Therefore, we have $9+9=18$ alternators, for a total of 31, which is satisfactory. On the other hand, if we have 12 truthtellers, that must mean we have 1 liar, meaning we have $9+14 = 23$ alternators. That count is 36, which is a contradiction. The answer is $\boxed{7}$.

 \textbf{Remark (Fake Solve)} This problem is broken in an interesting way that helps the test-taker. Since the true answers alternate for alternaters, you can still get the correct answer if you misinterpret the problem as "alternaters alternate their yes/no answers, but not their truth values". Suppose that $A$ alternaters answer yes-no-yes, and $A'$ alternaters answer no-yes-no. You still get $A=A'=9$. 

 In addition, you can even get the correct answer with the misinterpretation "alternaters have two types: alternaters of the same type give the same arbitrary answers to the three questions, and alternaters of different types give the opposite arbitrary answers to the three questions".

 It's also notable that the misinterpretation makes the problem harder, so that the solution actually relies on all the information. This suggests that the question-writer may have been mistaken but got lucky.

 


	\item (\textit{AMC 12 12A 2022}) How many ways are there to split the integers $1$ through $14$ into $7$ pairs such that in each pair, the greater number is at least $2$ times the lesser number?


 $\textbf{(A) } 108 \qquad \textbf{(B) } 120 \qquad \textbf{(C) } 126 \qquad \textbf{(D) } 132 \qquad \textbf{(E) } 144$

 \textbf{Answer: }E



\textbf{Solution 1 (Multiplication Principle)}

Clearly, the integers from $8$ through $14$ must be in different pairs, so are the integers from $1$ through $7.$ Note that $7$ must pair with $14.$

 We pair the numbers $1,2,3,4,5,6$ with the numbers $8,9,10,11,12,13$ systematically:

 
            \begin{itemize}
        \item $6$ can pair with either $12$ or $13.$
            \end{itemize}
            
            \begin{itemize}
        \item $5$ can pair with any of the three remaining numbers from $10,11,12,13.$
            \end{itemize}
            
            \begin{itemize}
        \item $1,2,3,4$ can pair with the other four remaining numbers from $8,9,10,11,12,13$ without restrictions.
            \end{itemize}
            Together, the answer is $2\cdot3\cdot4!=\boxed{\textbf{(E) } 144}.$




\textbf{Solution 2 (Multiplication Principle)}

As said in Solution 1, clearly, the integers from $8$ through $14$ must be in different pairs.

 We know that $8$ or $9$ can pair with any integer from $1$ to $4$, $10$ or $11$ can pair with any integer from $1$ to $5$, and $12$ or $13$ can pair with any integer from $1$ to $6$. Thus, $8$ will have $4$ choices to pair with, $9$ will then have $3$ choices to pair with ($9$ cannot pair with the same number as the one $8$ pairs with). $10$ cannot pair with the numbers $8$ and $9$ has paired with but can also now pair with $5$, so there are $3$ choices. $11$ cannot pair with $8$'s, $9$'s, or $10$'s paired numbers, so there will be $2$ choices for $11$. $12$ can pair with an integer from $1$ to $5$ that hasn't been paired with already, or it can pair with $6$. $13$ will only have one choice left, and $7$ must pair with $14$.

 So, the answer is $4\cdot3\cdot3\cdot2\cdot2\cdot1\cdot1=\boxed{\textbf{(E) } 144}.$




\textbf{Solution 3 (Generalization)}

The integers $x \in \{8, \ldots , 14 \}$ must each be the larger elements of a distinct pair.

 Assign partners in decreasing order for $x \in \{7, \dots, 1\}$:

 Note that $7$ must pair with $14$: $\mathbf{1} \textbf{ choice}$.

 For $5 \leq x \leq 7$, the choices are $\{2x, \dots, 14\} - \{ \text{previous choices}\}$. As $x$ decreases by 1, The minuend increases by 2 elements, and the subtrahend increases by 1 element, so the difference increases by 1, yielding $\mathbf{3!} \textbf{ combined choices}$.

 After assigning a partner to $5$, there are no invalid pairings for yet-unpaired numbers, so there are $\mathbf{4!} \textbf{ ways}$ to choose partners for $\{1,2,3,4\}$.

 The answer is $3! \cdot 4! = \boxed{\textbf{(E) } 144}$.

 In general, for $1,\ldots,2n$, the same logic yields answer: $\left\lfloor\dfrac{n}{2}\right\rfloor! \cdot \left\lceil\dfrac{n}{2}\right\rceil!$



	\item (\textit{AMC 12 12A 2022}) What is the product of all real numbers $x$ such that the distance on the number line between $\log_6x$ and $\log_69$ is twice the distance on the number line between $\log_610$ and $1$?


 $\textbf{(A) } 10 \qquad \textbf{(B) } 18 \qquad \textbf{(C) } 25 \qquad \textbf{(D) } 36 \qquad \textbf{(E) } 81$

 \textbf{Answer: }E



\textbf{Solution 1}

Let $a = 2 \cdot |\log_6 10 - 1| = |\log_6 9 - \log_6 x|  =  \left|\log_6 \frac{9}{x}\right|$.

 $\pm a = \log_6 \frac{9}{x} \implies 6^{\pm a} = b^{\pm 1} = \frac{9}{x}  \implies x = 9 \cdot b^{\pm 1}$
$9b^1 \cdot 9b^{-1} = \boxed{81}$.

 



\textbf{Solution 2}

First, notice that there must be two such numbers: one greater than $\log_69$ and one less than it. Furthermore, they both have to be the same distance away, namely $2(\log_610 - 1)$. Let these two numbers be $\log_6a$ and $\log_6b$. Because they are equidistant from $\log_69$, we have $\frac{\log_6a + \log_6b}{2} = \log_69$. Using log properties, this simplifies to $\log_6{\sqrt{ab}} = \log_69$. We then have $\sqrt{ab} = 9$, so $ab = \boxed{\textbf{(E) } 81}$.

 



\textbf{Solution 3 (Logarithmic Rules and Casework)}

In effect we must find all $x$ such that $\left|\log_6 9 - \log_6 x\right| = 2d$ where $d = \log_6 10 - 1$.

 Notice that by log rules
 \[d = \log_6 10 - 1 = \log_6 \frac{10}{6}\]Using log rules again,
 \[2d = 2\log_6 \frac{10}{6} = \log_6 \frac{25}{9}\]

 Now we proceed by casework for the distinct values of $x$.

 \textbf{Case 1} 
 \[\log_6 9 - \log_6 x_1 = 2d\]Subbing in for $2d$ and using log rules,
 \[\log_6 \frac{9}{x_1} = \log_6 \frac{25}{9}\]From this we may conclude that
 \[\frac{9}{x_1} = \frac{25}{9} \implies x_1 = \frac{81}{25}\]

 \textbf{Case 2} 
 \[\log_6 9 - \log_6 x_2 = -2d\]Subbing in for $-2d$ and using log rules,
 \[\log_6 \frac{9}{x_2} = \log_6 \frac{9}{25}\]From this we conclude that
 \[\frac{9}{x_2} = \frac{9}{25} \implies x_2 = 25\]

 Finding the product of the distinct values,$x_1x_2 = \boxed{\textbf{(E) } 81}$



	\item (\textit{AMC 12 12A 2022}) Let $M$ be the midpoint of $\overline{AB}$ in regular tetrahedron $ABCD$. What is $\cos(\angle CMD)$?


 $\textbf{(A) } \frac14 \qquad \textbf{(B) } \frac13 \qquad \textbf{(C) } \frac25 \qquad \textbf{(D) } \frac12 \qquad \textbf{(E) } \frac{\sqrt{3}}{2}$

 \textbf{Answer: }B



\textbf{Solution 1 (Right Triangles)}

Without loss of generality, let the edge-length of $ABCD$ be $2.$ It follows that $MC=MD=\sqrt3.$

 Let $O$ be the center of $\triangle ABD,$ so $\overline{CO}\perp\overline{MOD}.$ Note that $MO=\frac13 MD=\frac{\sqrt{3}}{3}.$

 In right $\triangle CMO,$ we have 
 \[\cos(\angle CMD)=\frac{MO}{MC}=\boxed{\textbf{(B) } \frac13}.\]~MRENTHUSIASM

 



\textbf{Solution 2 (Law of Cosines)}

Without loss of generality, let the edge-length of $ABCD$ be $2.$ It follows that $CM=DM=\sqrt3.$

 By the Law of Cosines, 
 \[\cos(\angle CMD) = \frac{CM^2 + DM^2 - CD^2}{2(CM)(DM)} = \boxed{\textbf{(B) } \frac13}.\]~jamesl123456

 



\textbf{Solution 3 (Double Angle Identities)}

As done above, let the edge-length equal $2$ (usually better than $1$ because we can avoid fractions when dropping altitudes). Notice that the triangle stated in the question has two side-lengths that are the altitudes of two equilateral triangles. By dropping the equilateral triangles' altitude and using $30^{\circ}$-$60^{\circ}$-$90^{\circ}$ properties, we find that the other two sides are equal to $\sqrt{3}$. Now by dropping the main triangle's altitude, we see it equals $\sqrt{2}$ from the Pythagorean Theorem. we can use the Double Angle Identities for Cosine. Doing this, we obtain 
 \[\cos(\angle CMD) = \frac{2}{3} - \frac13 = \boxed{\textbf{(B) } \frac13}.\]~Misclicked

 



\textbf{Solution 4 (Vector Methods)}

Without loss of generality, let tetrahedron $ABCD$ lie in three-dimensional space with vertices $A(-1, 1, 1)$, $B(1, 1, -1)$, $C(1, -1, 1)$, and $D(-1, -1, -1)$.  Let point $M$ be located at $(0, 1, 0)$.

 
 \[\vec{MD} = \langle -1, -2, -1 \rangle\] 
 \[\vec{MC} = \langle 1, -2, 1 \rangle\]

 We know that the dot product of two vectors equals the product of their magnitudes multiplied by the cosine of the angle between them:

 
 \[\vec{MD} \cdot \vec{MC} = (-1)(1) + (-2)(-2) + (-1)(1) = -1 + 4 - 1 = 2\]

 Compute the magnitudes:

 
 \[||\vec{MD}|| = \sqrt{(-1)^2 + (-2)^2 + (-1)^2} = \sqrt{1 + 4 + 1} = \sqrt{6}\] 
 \[||\vec{MC}|| = \sqrt{(1)^2 + (-2)^2 + (1)^2} = \sqrt{1 + 4 + 1} = \sqrt{6}\]

 Then:

 
 \[\vec{MD} \cdot \vec{MC} = ||\vec{MD}|| \cdot ||\vec{MC}|| \cdot \cos(\angle CMD)\]

 \[2 = \sqrt{6} \cdot \sqrt{6} \cdot \cos(\angle CMD) = 6 \cos(\angle CMD)\]

 Finally: 
 \[\cos(\angle CMD) = \frac{2}{6} = \boxed{\textbf{(B) } \frac13}.\]~TylerTrikowsky

 


	\item (\textit{AMC 12 12A 2022}) Let $\mathcal{R}$ be the region in the complex plane consisting of all complex numbers $z$ that can be written as the sum of complex numbers $z_1$ and $z_2$, where $z_1$ lies on the segment with endpoints $3$ and $4i$, and $z_2$ has magnitude at most $1$. What integer is closest to the area of $\mathcal{R}$?  


 $\textbf{(A) } 13 \qquad \textbf{(B) } 14 \qquad \textbf{(C) } 15 \qquad \textbf{(D) } 16 \qquad \textbf{(E) } 17$

 \textbf{Answer: }A



\textbf{Solution}


                    \begin{figure}[H]
                        \centering
                        \includegraphics[width=0.5\textwidth]{Solution/AMC_12_12A_2022_13_Solution_Figure_1.png}
                    \end{figure}
                    

 If $z$ is a complex number and $z = a + bi$, then the magnitude (length) of $z$ is $\sqrt{a^2 + b^2}$. Therefore, $z_1$ has a magnitude of 5. If $z_2$ has a magnitude of at most one, that means for each point on the segment given by $z_1$, the bounds of the region $\mathcal{R}$ could be at most 1 away. Alone the line, excluding the endpoints, a rectangle with a width of 2 and a length of 5, the magnitude, would be formed. At the endpoints, two semicircles will be formed with a radius of 1 for a total area of $\pi \approx 3$.Therefore, the total area is $5(2) + \pi \approx 10 + 3 = \boxed{\textbf{(A) } 13}$.

 


	\item (\textit{AMC 12 12A 2022}) What is the value of 
 \[(\log 5)^3+(\log 20)^3+(\log 8)(\log 0.25)\] where $\log$ denotes the base-ten logarithm?


 $\textbf{(A) } \frac{3}{2} \qquad \textbf{(B) } \frac{7}{4} \qquad \textbf{(C) } 2 \qquad \textbf{(D) } \frac{9}{4} \qquad \textbf{(E) } 3$

 \textbf{Answer: }C



\textbf{Solution 1}

Let $\text{log } 2 = x$. The expression then becomes 
 \[(1+x)^3+(1-x)^3+(3x)(-2x)=\boxed{2}.\]

 -bluelinfish

 



\textbf{Solution 1.1}

(Elaboration & motivation behind Sol. 1)

 Note that $5,20,8$, and $0.25$ are all products and quotients of exponents of $2$ and $5$, and the base of the logarithms is $10 = 2\times5$; this strongly hints at some sort of major simplification using the addition and subtraction rules of logarithms so we can convert all the different arguments of the logs into 1 common argument for easy algebra. 

 Note that we can write all of the following expressions in the following ways: 

 \begin{align*}\log5=\log\dfrac{10}2=\log10-\log2&=1-\log2\\ \log20=\log(2\cdot10)=\log2+\log10&=\log2+1\\\log8=\log\left(2^3\right)&=3\log2 \\\log0.25=\log\left(2^{-2}\right)&=-2\log2\end{align*}


 Thus, let $a=\log2$, and proceed as in solution 1. 

 



\textbf{Solution 2}

Using sum of cubes
 \[(\log 5)^{3}+(\log 20)^{3}\]
 \[= (\log 5 + \log 20)((\log 5)^{2}-(\log 5)(\log 20) + (\log 20)^{2})\]
 \[= 2((\log 5)^{2}-(\log 5)(2\log 2 + \log 5) + (2\log 2 + \log 5)^{2})\]Let x = $\log 5$ and y = $\log 2$, so $x+y=1$

 The entire expression becomes
 \[2(x^2-x(2y+x)+(2y+x)^2)-6y^2\]
 \[=2(x^2+2xy+4y^2-3y^2)\]
 \[=2(x+y)^2 = \boxed{2}\]




\textbf{Solution 3 (Estimates)}

We can estimate the solution. Using $\log(2) \approx 0.3, \log(20) = \log(10)+\log(2) = 1 + 0.3 \approx 1.3, \log(8) \approx 0.9$ and $\log(.25) = \log(1)-\log(4)= 0 - 0.6\approx -0.6,$ we have

 



\textbf{Solution 4(log bash)}

Using log properties, we combine the terms to make our expression equal to $\log {( (5^{\log^2{5}}) \cdot (20^{\log^2{20}}) \cdot 8 ^ {\log{\frac{1}{4}}} ) }$. By exponent properties, we separate the part with base $20$ to become $20^{\log^2{5}} \cdot 20^{\log^2{20}-\log^2{5}}$. Then, we substitute this into the original expression to get $\log {( (5^{\log^2{5}}) \cdot 20^{\log^2{5}} \cdot 20^{\log^2{20}-\log^2{5}} \cdot 8 ^ {\log{\frac{1}{4}}} ) } = \log {( (100^{\log^2{5}}) \cdot 20^{\log^2{20}-\log^2{5}} \cdot 8 ^ {\log{\frac{1}{4}}} ) }$. Because $100^{\log^2{5}} = 25^{\log{5}}$, and $\log^2{20}-\log^2{5} = (\log{20}+\log{5})(\log{20}-\log{5}) = \log{100}\cdot\log{4} = 2\log{4}$, this expression is equal to $\log {( 25 ^ {\log{5}} \cdot 400^{\log{4}} \cdot 8 ^ {\log{\frac{1}{4}}} ) }$. We perform the step with the base combining on $25$ and $400$ to get $25 ^ {\log{5}} \cdot 400^{\log{4}}  = 25 ^ {\log{5}-\log{4}} \cdot 10000^{\log{4}} = 25^{\log{\frac{5}{4}}}\cdot 256$. Putting this back into the whole equation gives $\log{( 25^{\log{\frac{5}{4}}}\cdot 256 \cdot 8^{\log{\frac{1}{4}}})}$. One last base merge remains - but $25\cdot 8$ isn't a power of 10. We can rectify this by converting $8^{\log{\frac{1}{4}}}$ to $(4^\frac{3}{2})^{\log{\frac{1}{4}}} = 4^{\log{ \frac{1}{8} }}$. Finally, we complete this arduous process by performing the base merge on $\log{( 25^{\log{\frac{5}{4}}}\cdot 256 \cdot 4^{\log{\frac{1}{8}}})}$. We get $25^{\log{\frac{5}{4}}} \cdot 4^{\log{\frac{1}{8}}} = 25^{\log{\frac{5}{4}}-\log{\frac{1}{8}}} \cdot 100^{\log{\frac{1}{8}}} = 25^{\log{10}} \cdot \frac{1}{64} = \frac{25}{64}$. Putting this back into that original equation one last time, we get $\log(256 \cdot \frac{25}{64}) = \log{100} = \boxed{2}$.~aop2014

 


	\item (\textit{AMC 12 12A 2022}) The roots of the polynomial $10x^3-39x^2+29x-6$ are the height, length, and width of a rectangular box (right rectangular prism). A new rectangular box is formed by lengthening each edge of the original box by $2$ units. What is the volume of the new box?


 $\textbf{(A) } \frac{24}{5} \qquad \textbf{(B) } \frac{42}{5} \qquad \textbf{(C) } \frac{81}{5} \qquad \textbf{(D) } 30 \qquad \textbf{(E) } 48$

 \textbf{Answer: }D



\textbf{Solution 1 (Vieta's Formulas)}

Let $a$, $b$, $c$ be the three roots of the polynomial. The lengthened prism's volume is 
 \[V = (a+2)(b+2)(c+2) = abc+2ac+2ab+2bc+4a+4b+4c+8 = abc + 2(ab+ac+bc) + 4(a+b+c) + 8.\]By Vieta's formulas, we know that a cubic polynomial $Ax^3+Bx^2+Cx+D$ with roots $a$, $b$, $c$ satisfies:
 \begin{alignat*}{8} a+b+c &= -\frac{B}{A} &&= \frac{39}{10}, \\ ab+ac+bc &= \hspace{2mm}\frac{C}{A} &&= \frac{29}{10}, \\ abc &= -\frac{D}{A} &&= \frac{6}{10}. \end{alignat*}We can substitute these into the expression, obtaining 
 \[V = \frac{6}{10} + 2\left(\frac{29}{10}\right) + 4\left(\frac{39}{10}\right) + 8 = \boxed{\textbf{(D) } 30}.\]

 - phuang1024

 



\textbf{Solution 2 (Guessing Roots)}

From the answer choices, we can assume the roots are rational numbers, and therefore this polynomial should be easily factorable.The coefficients of $x$ must multiply to $10$, so these coefficients must be $5,2,1$ or $10,1,$ in some order. We can try one at a time, and therefore write the factored form as follows: 
 \[(5x-p)(2x-q)(x-r).\]Note that $p, q, r$ have to multiply to $6$, so they must be either $3,2,1$ or $6,1,1$ in some order. Again, we can try one at a time in different positions and see if they multiply correctly.We try $(5x-2)(2x-1)(x-3)$ and multiply the $x-$terms, and sure enough they add up to $29$. You can try to add up the $x^2$ terms and they add up to $-39$. Therefore the roots are $\frac{2}{5}$, $\frac{1}{2}$ and $3$. Now if you add $2$ to each root, you get the volume is $\frac{12}{5} \cdot \frac{5}{2} \cdot 5 = 6 \cdot 5 = 30 = \boxed{\textbf{(D) } 30}$.

 



\textbf{Solution 3 (Rational Root Theorem Bash)}

We can find the roots of the cubic using the Rational Root Theorem, which tells us that the rational roots of the cubic must be in the form $\frac{p}{q}$, where $p$ is a factor of the constant $(-6)$ and $q$ is a factor of the leading coefficient $(10)$. Therefore, $p$ is $\pm (1, 2, 3, 6)$ and q is $\pm (1, 2, 5, 10).$

 Doing Synthetic Division, we find that $3$ is a root of the cubic:
 \[\begin{array}{c|rrrr}&10&-39&29&-6\\3&&30&-27&6\\\hline\\&10&-9&2&0\\\end{array}.\]

 Then, we have a quadratic $10x^2-9x+2.$ Using the Quadratic Formula, we can find the other two roots:
 \[x=\frac{9 \pm \sqrt{(-9)^2-4(10)(2)}}{2 \cdot 10},\]which simplifies to $x=\frac{1}{2}, \frac{2}{5}.$

 To find the new volume, we add $2$ to each of the roots we found:
 \[(3+2)\cdot\left(\frac{1}{2}+2\right)\cdot\left(\frac{2}{5}+2\right).\]Simplifying, we find that the new volume is $\boxed{\textbf{(D) } 30}.$

 -MathWizard09

 



\textbf{Solution 4}

Let $P(x) = 10x^3 - 39x^2 + 29x - 6$, and let $a, b, c$ be the roots of $P(x)$. The roots of $P(x-2)$ are then $a + 2, b + 2, c + 2,$ so the product of the roots of $P(x-2)$ is the area of the desired rectangular prism.

 $P(x-2)$ has leading coefficient $10$ and constant term $P(0-2) = P(-2) = 10(-2)^3 - 39(-2)^2 + 29(-2) - 6 = -300$. 

 Thus, by Vieta's Formulas, the product of the roots of $P(x-2)$ is $\frac{-(-300)}{10} = \boxed{\textbf{(D) } 30}$.

 -Orange\_Quail\_9

 



\textbf{Solution 5}

Let $P(x) = 10x^3 - 39x^2 + 29x - 6$. This can be factored m as $P(x) = 10(x-a)(x-b)(x-c)$, where $a$, $b$, and $c$ are the roots of $P(x)$. We want $V = (a+2)(b+2)(c+2)$.

 "Luckily" $P(-2) = 10(-2-a)(-2-b)(-2-c) = -10V$.  $P(-2) = -300$, giving $V = \boxed{\textbf{(D) } 30}$.  

 -Oxymoronic15

 (It's not just lucky. If $P(x)$ has roots $x \in \{r_i\}$, $Q(x) = P(x-2)$ has roots $y \in \{r_i+2\}$. By Vieta, the product $V$ of the roots is the negation of the constant term divided by the leading coefficient $10$ , which is $-Q(0) / 10$, which is $-P(0-2) / 10$. -oinava )

 



\textbf{Solution 6 (Very Nice Vieta Solve)}

If we let the roots of the cubic polynomial be $p,q,r$, then we can write what we want in terms of $p, q, r$. So, we have 
 \[(p+2)(q+2)(r+2) = pqr + 2pq + 2qr + 4q + 2pr + 4p + 4r + 8 = pqr + 2(pq + qr + pr) + 4(q + p + r) + 8.\] Now, we can just use Vieta's to finish up from here! We have $pqr = -(-\dfrac{6}{10}) = \dfrac{3}{5}$. Now, we can find the sum of the pairwise roots which is just $pq + qr + pr = \dfrac{29}{10}$ and then multiply by $2$ which gives $\dfrac{29}{5}$. Looking good so far! Now, we need to find the sum of the roots which is just $p + q + r = \dfrac{39}{10}$ and then multiply by $4$ which gives $\dfrac{78}{5}$. Almost there! Now, we just add $8$ to that sum, and then sum it all up to get $\boxed{\textbf{(D) } 30}$

 -jb2015007

 Note: This is the same as Solution 1.

 


	\item (\textit{AMC 12 12A 2022}) A triangular number is a positive integer that can be expressed in the form $t_n=1+2+3+\cdots+n$, for some positive integer $n$. The three smallest triangular numbers that are also perfect squares are $t_1=1=1^2, t_8=36=6^2,$ and $t_{49}=1225=35^2$. What is the sum of the digits of the fourth smallest triangular number that is also a perfect square?


 $\textbf{(A)} ~6 \qquad\textbf{(B)} ~9 \qquad\textbf{(C)} ~12 \qquad\textbf{(D)} ~18 \qquad\textbf{(E)} ~27$

 \textbf{Answer: }D



\textbf{Solution 1(Pell Equation)}

We have $t_n = \frac{n (n+1)}{2}$.If $t_n$ is a perfect square, then it can be written as$\frac{n (n+1)}{2} = k^2$,where $k$ is a positive integer.

 Thus, $n (n+1) = 2 k^2$. Rearranging, we get $(2n+1)^2-2(2k)^2=1$, a Pell equation (see \href{https://artofproblemsolving.com/wiki/index.php/Pell\_equation}{https://artofproblemsolving.com/wiki/index.php/Pell\_equation}). So $\frac{2n+1}{2k}$ must be a truncation of the continued fraction for $\sqrt{2}$:

 
 \begin{eqnarray*} 1+\frac12&=&\frac{2\cdot1+1}{2\cdot1}\\ 1+\frac1{2+\frac1{2+\frac12}}&=&\frac{2\cdot8+1}{2\cdot6}\\ 1+\frac1{2+\frac1{2+\frac1{2+\frac1{2+\frac12}}}}&=&\frac{2\cdot49+1}{2\cdot35}\\ 1+\frac1{2+\frac1{2+\frac1{2+\frac1{2+\frac1{2+\frac1{2+\frac12}}}}}}&=&\frac{2\cdot288+1}{2\cdot204} \end{eqnarray*}

 Therefore, $t_{288} = \frac{288\cdot289}{2} = 204^2 = 41616$, so the answer is $4+1+6+1+6=\boxed{\textbf{(D) 18}}$.

 
- Steven Chen (Professor Chen Education Palace, www.professorchenedu.com) 

 Edited by wzs26843545602Edited by dad

 



\textbf{Solution 1a (Alt. Pell)}

As in Solution 1, we find that $(2n+1)^2-8k^2=1$. Let $q=2n+1$; applying the difference of squares, we see that $(q-2\sqrt2k)(q+2\sqrt2k)=1$. For this Pell equation, note that $q+2\sqrt2k$ and $q-2\sqrt2k$ are both algebraic integers that are conjugates of each other. WLOG let $z=q-2\sqrt2k$ so that $\overline z=q+2\sqrt2k$; we know that $z\overline z=1$. It is clear that $\overline{ab}=\overline a\overline b$, and it immediately follows that $\left(z^n\right)\overline{\left(z^n\right)}=1$. In essence, for a solution $z$ to $z\overline z=1$, then any integer power of $z$ is also a solution. This is particularly useful in this problem: we are given that $n=1\implies q=3$ and $k=1$ is a solution, so $z=3-2\sqrt2$. From our result, we know that $z^2=\left(3-2\sqrt2\right)^2=17-2\sqrt2\cdot6$ is also a solution; this particular solution is $n=8$ and $k=6$, which is the second smallest desired number. Similarly, $z^3$ gives us the third-smallest desired number, so $z^4$ must give us our desired answer. We can calculate: $z^4=\left(17-2\sqrt2\cdot6\right)^2=289+288-2\cdot17\cdot6\cdot2\sqrt2=577-2\sqrt2\cdot204$. This lets $n=288$ and $k=204$; we thus know that $T_{288}=204^2=41616$, which has digit sum $\boxed{\textbf{(D) }18}$. QED. $\Box$




\textbf{Solution 2 (Bash)}

As mentioned above, $t_n = \frac{n (n+1)}{2}$. If $t_n$ is a perfect square, one of two things must occur when the fraction is split into a product. Either $\frac{n}{2}$ and $n+1$ must both be squares, or $n$ and $\frac{n+1}{2}$ must both be squares, and thus the search for the next perfect square triangular number can be narrowed down by testing values of $n$ that are close to or are perfect squares. After some work, we reach $n = 288$, $1$ less than $289$, and $t_{288} = \frac{288\cdot289}{2} = 144 \cdot 289 = 41616$. This product is a perfect square, and thus the sum of the digits of the fourth smallest perfect square triangular number is therefore $4+1+6+1+6=\boxed{\textbf{(D) 18}}$. 

 



\textbf{Solution 3}

According to the problem, we want to find integer $p$ such $\frac{n(n+1)}{2}=p^2$, after expanding, we have $n^2+n=2p^2, 4n^2+4n=8p^2, (2n+1)^2-8p^2=1$, we call $2n+1=q$, the equation becomes $q^2-8p^2=1$, obviously $(q,p)=(3,1)$ is the elementary solution for this pell equation, thus the forth smallest solution set $q_4+2\sqrt{2}p_4=(3+2\sqrt{2})^4=577+408\sqrt{2}$, which indicates $p=204, p^2=41616$ leads to $\boxed{18}$




\textbf{Solution 4}

If $n \choose 2$ is a square, then ${(2n-1)^2 \choose 2}$ is also a square. We can prove this quite simply:

 
 \[{(2n-1)^2 \choose 2}\]
 \[= \frac{(2n-1)^2 \cdot ((2n-1)^2 - 1)}{2}\]
 \[= \frac{(2n-1)^2 \cdot (2n \cdot (2n - 2))}{2}\]
 \[= (2n-1)^2 \cdot 4{n \choose 2}.\]

 Therefore, ${(2 \cdot 9 - 1)^2 \choose 2}$ is a square. Note that $T_n = {n+1 \choose 2}$. We can easily check all smaller possibilities using a bit of casework, and they don't work. Our solution is thus ${289 \choose 2} = 204^2 = 41616$, and so the answer is $\boxed{18}$.

 



\textbf{Solution 5}

We want to find integer $n_i$ and $m_i$ such that $t_{n_i} =\frac{n_i (n_i + 1)}{2}=m_i^2, n_0 = 0.$

 We use the formula $\sqrt{n_{i+1}} = \sqrt{2n_i} + \sqrt{n_i + 1}$ and get

 
 \[n_1 = ( \sqrt{2n_0} + \sqrt{n_0 + 1})^2 = (0+1)^2 = 1,\] 
 \[n_2= ( \sqrt{2n_1} + \sqrt{n_1 + 1})^2 = ( \sqrt{2} + \sqrt{1 + 1})^2 = 8,\]
 \[n_3= ( \sqrt{2n_2} + \sqrt{n_2 + 1})^2 = ( \sqrt{16} + \sqrt{8 + 1})^2 = 49,\]
 \[n_4 = ( \sqrt{2n_3} + \sqrt{n_3 + 1})^2 = ( \sqrt{98} + \sqrt{49 + 1})^2 =  ((7 + 5)\sqrt{2})^2 = 288,\]
 \[n_5 = ( \sqrt{2n_4} + \sqrt{n_4 + 1})^2 = ( \sqrt{576} + \sqrt{289})^2 =  (24 + 17)^2 = 1681,\]
 \[n_6 = ( \sqrt{2n_5} + \sqrt{n_5 + 1})^2 = ( 41\sqrt{2} + \sqrt{1682})^2 =  ((41 + 29)\sqrt{2})^2 = 9800,...\]Therefore, $t_{n_4} = t_{288} = \frac{288\cdot289}{2} = 41616 \implies 4+1+6+1+6=\boxed{\textbf{(D) 18}}$

            \begin{itemize}
        \item \href{https://artofproblemsolving.com/wiki/index.php?title=Pell\%27s\_equation\_(simple\_solutions)}{Pell's equation (simple solutions)}
            \end{itemize}
            \textbf{vladimir.shelomovskii@gmail.com, vvsss}

 



\textbf{Solution 6(Fast and Simple)}

This is an easy one to tackle by just noticing a pattern. Write out all the perfect squares before 300, and then notice that 4 and 9 are what are the factors of the 8th term. 25 and 49 are factors of the 49th term. 4 multiplied by 2 is one less then 9. 25 multiplied by 2 is one more than 49. Well, then we keep on writing out the perfect squares and then hunt for any that can be multiplied by two and then has an absolute value of 1 valued difference. 144 and 289 satisfy this. 288th term would give us these factors. By the way, this isn't  very rigorous. 

 No formula needed. 

 emilyyunhanq@gmail.com

 


	\item (\textit{AMC 12 12A 2022}) Suppose $a$ is a real number such that the equation 
 \[a\cdot(\sin{x}+\sin{(2x)}) = \sin{(3x)}\]
has more than one solution in the interval $(0, \pi)$. The set of all such $a$ that can be written
in the form 
 \[(p,q) \cup (q,r),\]
where $p, q,$ and $r$ are real numbers with $p < q< r$. What is $p+q+r$?


 $\textbf{(A) } {-}4 \qquad \textbf{(B) } {-}1 \qquad \textbf{(C) } 0 \qquad \textbf{(D) } 1 \qquad \textbf{(E) } 4$

 \textbf{Answer: }A



\textbf{Solution 1}

We are given that $a\cdot(\sin{x}+\sin{(2x)})=\sin{(3x)}$

 Using the sine double angle formula combine with the fact that $\sin{(3x)} = \sin{x}\cdot(4\cos^2{x}-1)$, which can be derived using sine angle addition with $\sin{(2x + x)}$, we have 
 \[a\cdot(\sin{x}+2\sin{x}\cos{x})=\sin{x}\cdot(4\cos^2{x}-1)\]Since $\sin{x} \ne 0$ as it is on the open interval $(0, \pi)$, we can divide out $\sin{x}$ from both sides, leaving us with 
 \[a\cdot(1+2\cos{x})=4\cos^2{x}-1\]Now, distributing $a$ and rearranging, we achieve the equation 
 \[4\cos^2{x} - 2a\cos{x} - (1+a) = 0\] which is a quadratic in $\cos{x}$.

 Applying the quadratic formula to solve for $\cos{x}$, we get 
 \[\cos{x} =\frac{2a\pm\sqrt{4a^2+4(4)(1+a)}}{8}\] and expanding the terms under the radical, we get 
 \[\cos{x} =\frac{2a\pm\sqrt{4a^2+16a+16}}{8}\]Factoring, since $4a^2+16a+16 = (2a+4)^2$, we can simplify our expression even further to 
 \[\cos{x} =\frac{a\pm(a+2)}{4}\]

 Now, solving for our two solutions,  $\cos{x} = -\frac{1}{2}$ and $\cos{x} = \frac{a+1}{2}$.

 Since $\cos{x} = -\frac{1}{2}$ yields a solution that is valid for all $a$, that being $x = \frac{2\pi}{3}$, we must now solve for the case where $\frac{a+1}{2}$ yields a valid value.

 As $x\in (0, \pi)$, $\cos{x}\in (-1, 1)$, and therefore $\frac{a+1}{2}\in (-1, 1)$, and $a\in(-3,1)$.

 There is one more case we must consider inside this interval though, the case where $\frac{a+1}{2} = -\frac{1}{2}$, as this would lead to a double root for $\cos{x}$, yielding only one valid solution for $x$. Solving for this case, $a \ne -2$.

 Therefore, combining this fact with our solution interval, $a\in(-3, -2) \cup (-2, 1)$, so the answer is $-3-2+1 = \boxed{\textbf{(A) } {-}4}$.

 - DavidHovey

 



\textbf{Solution 2}

We can optimize from the step from 
 \[a\cdot(1+2\cos{x})=4\cos^2{x}-1\] in solution 1 by writing

 
 \[a = \frac{4\cos^2{x}-1}{1+2\cos{x}} = 2\cos x - 1\]

 and then get
 \[\cos x = \frac{a+1}{2}.\]

 Now, solving for our two solutions,  $\cos{x} = -\frac{1}{2}$ and $\cos{x} = \frac{a+1}{2}$.

 Since $\cos{x} = -\frac{1}{2}$ yields a solution that is valid for all $a$, that being $x = \frac{2\pi}{3}$, we must now solve for the case where $\frac{a+1}{2}$ yields a valid value.

 As $x\in (0, \pi)$, $\cos{x}\in (-1, 1)$, and therefore $\frac{a+1}{2}\in (-1, 1)$, and $a\in(-3,1)$.

 There is one more case we must consider inside this interval though, the case where $\frac{a+1}{2} = -\frac{1}{2}$, as this would lead to a double root for $\cos{x}$, yielding only one valid solution for $x$. Solving for this case, $a \ne -2$.

 Therefore, combining this fact with our solution interval, $a\in(-3, -2) \cup (-2, 1)$, so the answer is $-3-2+1 = \boxed{\textbf{(A) } {-}4}$.

 - Dan

 Remark: If you get to $a \in (-3, 1)$, then you can observe from the answer choices that $q$ must be $-2, 1,$ or something greater than $1$. Then you can guess that $q=-2$ since $q$ can't be $1$.

 



\textbf{Solution 3}

Use the sum to product formula to obtain $2a\cdot\sin{\frac{3x}{2}}\cos{\frac{x}{2}}=\sin{3x}$. Use the double angle formula on the RHS to obtain $a\cdot\sin{\frac{3x}{2}}\cos{\frac{x}{2}}=\sin{\frac{3x}{2}}\cos{\frac{3x}{2}}$. From here, it is obvious that $x=\frac{2\pi}{3}$ is always a solution, and thus we divide by $\sin{\frac{3x}{2}}$ to get 
 \[a\cdot\cos{\frac{x}{2}}=\cos{\frac{3x}{2}}\] We wish to find all $a$ such that there is at least one more solution to this equation distinct from $x=\frac{2\pi}{3}$. Letting $y=\cos{\frac{x}{2}}$, and noting that $\cos{\frac{3x}{2}}=4y^3-3y$, we can rearrange our equation to $4y^3=y(3+a)$ The smallest value $x$ where $y=0$ is $\pi$, which is not in our domain so we divide by $y$ to obtain $4y^2=a+3$. By the trivial inequality, $a+3\ge{0}$. Furthermore, $y\neq{0}$, so $a+3>0$. Also, if $a=-2$, then the solution to this equation would be shared with $x=\frac{2\pi}{3}$, so there would only be one distinct solution. Finally, because $y\le{1}$ due to the restrictions of a sine wave, and that $y\neq{1}$ due to the restrictions on $x$, we have $-3<a<1$ with $a\neq{-2}$. Thus, $p=-3,q=-2, r=1$, so our final answer is $-3+(-2)+1=\boxed{\textbf{(A) } {-}4}$. 

 


	\item (\textit{AMC 12 12A 2022}) Let $T_k$ be the transformation of the coordinate plane that first rotates the plane $k$ degrees counterclockwise around the origin and then reflects the plane across the $y$-axis. What is the least positive integer $n$ such that performing the sequence of transformations $T_1, T_2, T_3, \dots, T_n$ returns the point $(1,0)$ back to itself?


 $\textbf{(A) } 359 \qquad \textbf{(B) } 360\qquad \textbf{(C) } 719 \qquad \textbf{(D) } 720 \qquad \textbf{(E) } 721$

 \textbf{Answer: }A



\textbf{Solution 1 (Polar Coords)}

Let $P=(r,\theta)$ be a point in polar coordinates, where $\theta$ is in degrees.

 Rotating $P$ by $k^{\circ}$ counterclockwise around the origin gives the transformation $(r,\theta)\rightarrow(r,\theta+k^{\circ}).$ Reflecting $P$ across the $y$-axis gives the transformation $(r,\theta)\rightarrow(r,180^{\circ}-\theta).$ Note that 
 \begin{align*} T_k(P)&=(r,180^{\circ}-\theta-k^{\circ}), \\ T_{k+1}(T_k(P)) &= (r,\theta -1^{\circ}). \end{align*}
We start with $(1,0^{\circ})$ in polar coordinates. For the sequence of transformations $T_1, T_2, T_3, \cdots, T_k,$ it follows that

 
            \begin{itemize}
        \item After $T_1,$ we have $(1,179^{\circ}).$
            \end{itemize}
            
            \begin{itemize}
        \item After $T_2,$ we have $(1,-1^{\circ}).$
            \end{itemize}
            
            \begin{itemize}
        \item After $T_3,$ we have $(1,178^{\circ}).$
            \end{itemize}
            
            \begin{itemize}
        \item After $T_4,$ we have $(1,-2^{\circ}).$
            \end{itemize}
            
            \begin{itemize}
        \item After $T_5,$ we have $(1,177^{\circ}).$
            \end{itemize}
            
            \begin{itemize}
        \item After $T_6,$ we have $(1,-3^{\circ}).$
            \end{itemize}
            
            \begin{itemize}
        \item ...
            \end{itemize}
            
            \begin{itemize}
        \item After $T_{2k-1},$ we have $(1,180^{\circ}-k^{\circ}).$
            \end{itemize}
            
            \begin{itemize}
        \item After $T_{2k},$ we have $(1,-k^{\circ}).$
            \end{itemize}
            The least such positive integer $k$ is $180.$ Therefore, the least such positive integer $n$ is $2k-1=\boxed{\textbf{(A) } 359}.$




\textbf{Solution 2}

Note that since we're reflecting across the $y$-axis, if the point ever makes it to $(-1,0)$ then it will flip back to the original point. Note that after $T_1$ the point will be $1$ degree clockwise from the negative $x$-axis. Applying $T_2$ will rotate it to be $1$ degree counterclockwise from the negative $x$-axis, and then flip it so that it is $1$ degree clockwise from the positive $x$-axis. Therefore, after every $2$ transformations, the point rotates $1$ degree clockwise. To rotate it so that it will rotate $179$ degrees clockwise will require $179 \cdot 2 = 358$ transformations. Then finally on the last transformation, it will rotate on to $(-1,0)$ and then flip back to its original position. Therefore, the answer is $358+1 = 359 = \boxed{\textbf{(A) } 359}$.

 



\textbf{Solution 3}

So, you don't know polar coordinates (Solution 1) nor came up with Solution 2's logic, then, here we go!

 The easy solution is to notice that, without reflections and taking \( k=1 \), there are 360 turns that need to be made.

 Now take \( k=2 \). We see that there are 180 rotations, so the number of rotations is just \( 360/k \).

 Now include the reflections over the \( y \) axis.

 Say that each time we reflect a point over the \( y \) axis, we draw a line.

 Each time we reflect over the \( y\) axis, we see more and more lines that generate. For \( k=1 \), there are quite obviously 360 of these lines.

 However, one can see that this works for every single case, because when we reflect it over the \( y\) axis and make a \( k^\circ \) rotation, we fill up all the missed gaps until we reach the end of a sequence - I prove this on section 1.1.

 But wait! We overcounted by one line! Notice how I said there were 360 lines, but the line \(x=1\) intersects \( y=0\) is already given to you, and that is when \(T_k = 0\), which is the point \( (1, 0) \)! So our final answer is \( 360-1 = \) $\boxed{\textbf{(A) } 359}.$

 \textbf{Proof} A way you can prove such an assumption is by taking ideas like addition. We know that eventually there exists a number \( N \) such that \( l \times N = m \), where \( m \) has a units digit of \( l \). Utilizing the same approach, we can find that this actually works.

 



\textbf{Solution 4}

In degrees:

 Starting with $n=0$, the sequence goes ${0}\rightarrow {179}\rightarrow {359}\rightarrow {178}\rightarrow {358}\rightarrow {177}\rightarrow {357}\rightarrow\cdots.$

 We see that it takes $2$ steps to downgrade the point by $1^{\circ}$. Since the $1$st point in the sequence is ${179}$, the answer is $1+2(179)=\boxed{\textbf{(A) } 359}.$




\textbf{Solution 5 (Simple)}

We can consider the rotations and reflections separately. For the rotations, each rotation turns it by the next natural number. Thus the total number of degrees turned would be a triangle number. We test the smallest number, $359$ first, and we get that it turns $\frac{(1+359)359}{2} = 180n$, where $n$ is an integer. Thus, the point would be rotated to $(-1,0)$. We may be tempted to dismiss this option but we haven't considered the reflections. Each reflection acts as a $180^{\circ}$ rotation, so every two reflections cancel. However, $359$ is odd so we have to reflect $(-1,0)$, taking us to $(1,0)$, which is what we want. Thus we get $\boxed{\textbf{(A) } 359}$.

 



\textbf{Solution 6 (Complex)}

Rotations and reflections in 2D are very nice to describe using complex numbers in polar form. Reflections and rotations also don't affect the length (modulus or absolute value) of a complex number. This motivates us to set $z = \exp (i\theta)$, starting with $z=1+0i=\exp(i0)$.Then, we can describe rotations and reflections about the $y$-axis using these formulas:
 \[\text{Rotation}(\exp(i\theta), k): \exp(i\theta) \mapsto \exp(i(\theta + k))\]
 \[\text{Reflection}(\exp(i\theta)): \exp(i\theta) \mapsto \exp i (180^\circ - \theta) = -\exp(i(-\theta))\]

 If we apply two successive iterations, we see a simplification:

 
 \[T_k(\exp(i\theta)) = -\exp(-i(\theta+k))\]

 \begin{aligned}T\_{k+1}(T\_k(\exp(i\theta))) &= T\_{k+1}(-\exp(i(-\theta-k))\\&= --\exp(-i(-\theta-k+(k+1))\\&= \exp i(\theta-1)\end{aligned}

 We also can calculate $T_1(\exp i0) = \exp 179^\circ$. Thus, the point $(1,0)$ gets sent back to when all double iterations after $1$ cancel $179^\circ$. $179-\left(\frac{n-1}{2}\right)=0$ so $n = 1 + 2 \cdot 179 = \boxed{\textbf{(A) }359}$.

 



\textbf{Solution 7 (Matrices)}

Similar to the above solution, one can use rotation matrices to compute it.

 (i.e. $\begin{bmatrix} cos(\theta) & -sin(\theta) \\ sin(\theta) & cos(\theta) \end{bmatrix}$)

 For the reflection, use the reflection matrix (i.e. $\begin{bmatrix} -1 & 0 \\ 0 & 1 \end{bmatrix}$)

 Then proceed like Solution 5.

 



\textbf{Solution 8 (Very simple and quick)}

Let's test a few small cases for $T_k$. When $k=1$, the point will move $(1+180)^{\circ}$. When $k=2$, the point will move $(2+180)^{\circ}$. We can see that $T_k$ will result in the point moving $(k+180)^{\circ}$ from its previous location. Since we want the point to return to itself, the total number of degrees it has rotated must be a multiple of $360^{\circ}$. This means $T_1 + T_2 +T_3 + ... T_k \equiv 0 \pmod{360}$. Using the formula for the sum of the first $k$ integers, this can be expressed as 
 \[180k + \frac{k(k+1)}{2} \equiv 0 \pmod{360}\] Testing the answer choices gives $\boxed{\textbf{(A) } 359}.$



	\item (\textit{AMC 12 12A 2022}) Suppose that $13$ cards numbered $1, 2, 3, \ldots, 13$ are arranged in a row. The task is to pick them up in numerically increasing order, working repeatedly from left to right. In the example below, cards $1, 2, 3$ are picked up on the first pass, $4$ and $5$ on the second pass, $6$ on the third pass, $7, 8, 9, 10$ on the fourth pass, and $11, 12, 13$ on the fifth pass. For how many of the $13!$ possible orderings of the cards will the $13$ cards be picked up in exactly two passes?


 
                    \begin{figure}[H]
                        \centering
                        \includegraphics[width=0.5\textwidth]{Problem/AMC_12_12A_2022_19_Problem_Figure_1.png}
                    \end{figure}
                    
$\textbf{(A) } 4082 \qquad \textbf{(B) } 4095 \qquad \textbf{(C) } 4096 \qquad \textbf{(D) } 8178 \qquad \textbf{(E) } 8191$

 \textbf{Answer: }D



\textbf{Solution 1 (Casework)}

For $1\leq k\leq 12,$ suppose that cards $1, 2, \ldots, k$ are picked up on the first pass. It follows that cards $k+1,k+2,\ldots,13$ are picked up on the second pass.

 \textbf{Once we pick the spots for the cards on the first pass, there is only one way to arrange all $\boldsymbol{13}$ cards.} 

 For each value of $k,$ there are $\binom{13}{k}-1$ ways to pick the $k$ spots for the cards on the first pass: We exclude the arrangement
                    \begin{figure}[H]
                        \centering
                        \includegraphics[width=0.5\textwidth]{Solution/AMC_12_12A_2022_19_Solution_Figure_1.png}
                    \end{figure}
                    in which the cards are arranged such that the first pass consists of all $13$ cards.

 Therefore, the answer is 
 \[\sum_{k=1}^{12}\left[\binom{13}{k}-1\right] = \left[\sum_{k=1}^{12}\binom{13}{k}\right]-12 = \left[\sum_{k=0}^{13}\binom{13}{k}\right]-14 = 2^{13} - 14 = \boxed{\textbf{(D) } 8178}.\]




\textbf{Solution 2 (Casework)}

Since the $13$ cards are picked up in two passes, the first pass must pick up the first $n$ cards and the second pass must pick up the remaining cards $m$ through $13$. Also note that if $m$, which is the card that is numbered one more than $n$, is placed before $n$, then $m$ will not be picked up on the first pass since cards are picked up in order. Therefore we desire $m$ to be placed before $n$ to create a second pass, and that after the first pass, the numbers $m$ through $13$ are lined up in order from least to greatest.

 To construct this, $n$ cannot go in the $n$th position because all cards $1$ to $n-1$ will have to precede it and there will be no room for $m$. Therefore $n$ must be in slots $n+1$ to $13$.Let's do casework on which slot $n$ goes into to get a general idea for how the problem works.

 \textbf{Case 1:} With $n$ in spot $n+1$, there are $n$ available slots before $n$, and there are $n-1$ cards preceding $n$. Therefore the number of ways to reserve these slots for the $n-1$ cards is $\binom{n}{n-1}$. Then there is only $1$ way to order these cards (since we want them in increasing order). Then card $m$ goes into whatever slot is remaining, and the $13-m$ cards are ordered in increasing order after slot $n+1$, giving only $1$ way. Therefore in this case there are $\binom{n}{n-1}$ possibilities.

 \textbf{Case 2:} With $n$ in spot $n+2$, there are $n+1$ available slots before $n$, and there are $n-1$ cards preceding $n$. Therefore the number of ways to reserve slots for these cards are $\binom{n+1}{n-1}$. Then there is one way to order these cards. Then cards $m$ and $m+1$ must go in the remaining two slots, and there is only one way to order them since they must be in increasing order. Finally, cards $m+2$ to $13$ will be ordered in increasing order after slot $n+1$, which yields $1$ way. Therefore, this case has $\binom{n+1}{n-1}$ possibilities.

 I think we can see a general pattern now. With $n$ in slot $x$, there are $x-1$ slots to distribute to the previous $n-1$ cards, which can be done in $\binom{x-1}{n-1}$ ways. Then the remaining cards fill in in just $1$ way. Since the cases of $n$ start in slot $n+1$ and end in slot $13$, this sum amounts to:
 \[\binom{n}{n-1}+\binom{n+1}{n-1}+\binom{n+2}{n-1} + \cdots + \binom{12}{n-1}\] for any $n$.

 Hmmm ... where have we seen this before?

 We use wishful thinking to add a term of $\binom{n-1}{n-1}$:
 \[\binom{n-1}{n-1}+\binom{n}{n-1}+\binom{n+1}{n-1}+\binom{n+2}{n-1} + \cdots + \binom{12}{n-1}\]

 This is just the hockey stick identity! Applying it, this expression is equal to $\binom{13}{n}$. However, we added an extra term, so subtracting it off, the total number of ways to order the $13$ cards for any $n$ is 
 \[\binom{13}{n}-1\]

 Finally, to calculate the total for all $n$, we sum from $n=0$ to $13$. This yields us:

 
 \[\sum_{n=0}^{13} \binom{13}{n}-1 \implies \sum_{n=0}^{13} \binom{13}{n} - \sum_{n=0}^{13} 1\]
 \[\implies 2^{13} - 14 = 8192 - 14 = 8178 = \boxed{\textbf{(D) } 8178}.\]




\textbf{Solution 3 (Recursion)}

To solve this problem, we can use recursion on $n$. Let $A_n$ be the number of arrangements for $n$ numbers. Now, let's look at how these arrangements are related to $A_{n-1}$ by case work on the first number $a_1$.

 If $a_1 = 1$, the remaining $n-1$ numbers from $2$ to $n$ are arranged in the same way just like number 1 to $n-1$ in the case of $n-1$ numbers. So there are $A_{n-1}$ arrangements.

 If $a_1 = 2$, then we need to choose 1 position from position 2 to $n-1$ to put 1, and all remaining numbers must be arranged in increasing order, so there are $\binom{n-1}{1}$ such arrangements.

 If $a_1 = k$, then we need to choose $k-1$ positions from position 2 to $n-1$ to put $1, 2,\cdots k-1$, and all remaining numbers must be arranged in increasing order, so there are $\binom{n-1}{k-1}$ such arrangements.

 So we can write 
 \[A_n = A_{n-1} + \binom{n-1}{1} + \binom{n-1}{2} + \cdots + \binom{n-1}{n-1}\]which can be simplified to
 \[A_n = A_{n-1} + 2^{n-1} - 1\] We can solve this recursive sequence by summing up $n-1$ lines of the recursive formula
 \[A_n - A_{n-1} = 2^{n-1} - 1\]
 \[A_{n-1} - A_{n-2} = 2^{n-2} - 1\]
 \[\cdots\]
 \[A_2 - A_{1} = 2^{1} - 1\]to get
 \[A_n - A_1 = \sum_{k=1}^{n-1} (2^k - 1) = 2^n - 2 - (n-1) = 2^n - n - 1\]since $A_1=0$, we have
 \[A_n = 2^n - n - 1\]and $A_{13} = 2^{13} - 14 = \boxed{\textbf{(D) } 8178}$.

 
-- Dan Li

 



\textbf{Solution 4 (Engineer's Induction)}

When we have $3$ cards arranged in a row, after listing out all possible arrangements, we see that we have $4$ arrangements: $(1, 3, 2), (2, 1, 3), (2, 3, 1),$ and $(3, 1, 2)$. When we have $4$ cards, we find $11$ possible arrangements: $(1, 2, 4, 3), (1, 3, 2, 4), (1, 3, 4, 2), (1, 4, 2, 3), (2, 1, 3, 4), (2, 3, 1, 4), (2, 3, 4, 1), (3, 1, 2, 4), (3, 1, 4, 2), (3, 4, 1, 2),$ and $(4, 1, 2, 3).$ Hence, we recognize the pattern that for $n$ cards, we have $2^n - n - 1$ valid arrangements, so our answer is $2^{13} - 13 - 1 = \boxed{\textbf{(D) } 8178}.$ Note that this also works for $2$ cards being arranged in a row. $(2,1)$ Here, only one works, and following the formula, it checks out.

 



\textbf{Solution 5 (Subsets Bijection)}

Notice that for each card "position", we can choose for it to be picked up on the first or second pass, for a total of $2^{13}$ options. However, if all of the cards selected to be picked up first are before all of the cards to be picked up second, then this means that the list is in consecutive ascending order (and thus all cards will be picked up on the first pass instead). This can happen in 14 ways, so our answer is $2^{13}-14=\boxed{\textbf{(D) } 8178}$.

 



\textbf{Solution 6 (10 Seconds Solution)}

To satisfy the condition the question provided, we can arrange the numbers between 1 and 13 into 3 groups. Though writing out the groups is not necessary, below is a basic 'format'. (\# of elements in 1st group, \# of elements in the 2nd group, \# of elements in the 3rd group)(1,1,11)=(1,2,10)=(1,3,9)=(1,4,8)=(1,5,7)=(1,6,6)=(2,2,9)=(2,3,8)=(2,4,7)=(2,3,8)=(2,4,7)=(2,5,6)=(3,3,7)=(3,4,6)=(3,5,5)

 Then to arrange these groups we need to multiply the final value by 3!. Hence, we know that the answer should be a multiple of 3!, which is 6. Since D is the only answer choice that is a multiple of 6, the answer is $\boxed{D}$.~minor edit(Marshall\_Huang)

 



\textbf{Solution 7 (Simpler Casework)}

Let $x$ be the largest card picked up on the first pass. Then $1, 2, \ldots, x$ are picked up first and $x+1, \ldots, 13$ second, so define
 \[A = \{1, 2, \ldots, x\}, \qquad B = \{x+1, \ldots, 13\}.\]Both $A$ and $B$ must be non-empty such that we have $2$ passes, so $1 \le x \le 12$.

 
Replace each card by its label, $A$ or $B$, giving a row of $x$ $A$'s and $13-x$ $B$'s. For each such row of labels, there is exactly $1$ way to assign the numbers back onto the cards: inherently, the $A$-positions must read $1, 2, \ldots, x$ from left to right and the $B$-positions must read $x+1, \ldots, 13$ from left to right. So the arrangements we want correspond one-to-one with the label rows, and the number of label rows is
 \[\frac{13!}{x!\,(13-x)!} = \binom{13}{x},\]choosing which $x$ of the $13$ positions get an $A$. However, there is an arrangement with all the $A$ cards in a row followed by all the $B$ cards. Since order is fixed to be increasing, such a row reconstructs to $1, 2, \ldots, 13$, which takes $1$ pass instead of $2$, meaning we must subtract one case per value of $x$, or $12$.

 
Since arrangements of $A$ and $B$ between different values of $x$ are disjoint, we can sum without worrying about overcounting:
 \[\binom{13}{1} + \binom{13}{2} + \cdots + \binom{13}{12} - 12.\]

 The sum looks awfully similar to the Pascal's Triangle row-sum identity:
 \[\sum_{k=0}^{13}\binom{13}{k} = \binom{13}{0} + \binom{13}{1} + \ldots + \binom{13}{13} = 2^{13}\]

 With a little bit of wishful thinking and rearranging, we get
 \begin{align*} \binom{13}{0} + \binom{13}{1} + \cdots + \binom{13}{13} - \binom{13}{0} - \binom{13}{13} - 12 &= 2^{13} - \binom{13}{0} - \binom{13}{13} - 12 \\ \binom{13}{1} + \binom{13}{2} + \cdots + \binom{13}{12} - 12 &= 2^{13} - \binom{13}{0} - \binom{13}{13} - 12 \\ &= 8192 - 1 - 1 - 12 \\ &= \boxed{\textbf{(D) } 8178}. \end{align*}




	\item (\textit{AMC 12 12A 2022}) Isosceles trapezoid $ABCD$ has parallel sides $\overline{AD}$ and $\overline{BC},$ with $BC < AD$ and $AB = CD.$ There is a point $P$ in the plane such that $PA=1, PB=2, PC=3,$ and $PD=4.$ What is $\frac{BC}{AD}?$
$\textbf{(A) }\frac{1}{4}\qquad\textbf{(B) }\frac{1}{3}\qquad\textbf{(C) }\frac{1}{2}\qquad\textbf{(D) }\frac{2}{3}\qquad\textbf{(E) }\frac{3}{4}$

 \textbf{Answer: }B



\textbf{Solution 1 (Reflections + Ptolemy's Theorem)}

Consider the reflection 

 


	\item (\textit{AMC 12 12A 2022}) Let 
 \[P(x) = x^{2022} + x^{1011} + 1.\] Which of the following polynomials is a factor of $P(x)$?


 $\textbf{(A)} \, x^2 -x + 1 \qquad\textbf{(B)} \, x^2 + x + 1 \qquad\textbf{(C)} \, x^4 + 1 \qquad\textbf{(D)} \, x^6 - x^3 + 1  \qquad\textbf{(E)} \, x^6 + x^3 + 1$

 \textbf{Answer: }E



\textbf{Solution 1}

$P(x) = x^{2022} + x^{1011} + 1$ is equal to $\frac{x^{3033}-1}{x^{1011}-1}$ by difference of powers.

 Therefore, the answer is a polynomial that divides $x^{3033}-1$ but not $x^{1011}-1$.

 Note that any polynomial $x^m-1$ divides $x^n-1$ if and only if $m$ is a factor of $n$.

 The prime factorizations of $1011$ and $3033$ are $3*337$ and $3^2*337$, respectively.

 Hence, $x^9-1$ is a divisor of $x^{3033}-1$ but not $x^{1011}-1$.

 By difference of powers, $x^9-1=(x^3-1)(x^6+x^3+1)$.Therefore, the answer is $\boxed{E}$.

 



\textbf{Solution 2}

We simply test roots for each, as $2022,1011$ are multiples of three, we need to make sure the roots are in the form of $e^{i\frac{k\pi}{9}}$, so we only have to look at $D,E$.

 If we look at choice $E$, $x=e^{i\frac{\pm2\pi}{9}}$ which works perfectly, the answer is just $E$.

 



\textbf{Solution 3}

Let $x^{1011} = u$, now we can rewrite our polynomial as $u^2+u+1$. Using the quadratic formula to solve for the roots of this polynomial, we have 
 \[x^{1011} = \frac{-1\pm i\sqrt{3}}{2}\] Looking at our answer choices, we want to find a polynomial whose roots satisfy this expression. Since the expression $x^6+x^3+1$ is in a similar form to our original polynomial, except with $x^3$ in place of $x^{1011}$, this would be a good place to start. Solving for the roots of $x^3$ in a similar fashion, 
 \[x^3= \frac{-1\pm i\sqrt{3}}{2}\] for the solution we are testing. Now notice that we can rewrite the roots of $x^3$ as 
 \[x^3 = \operatorname{cis}{\frac{2\pi}{3}}, \operatorname{cis}{\frac{4\pi}{3}}\] Both of which are third roots of unity. We want to now check if this value of $x^3$ satisfies $x^{1011} = \frac{-1\pm i\sqrt{3}}{2}$. Notice that $x^{1011} = (x^{3})^{112\cdot3}\cdot x^3$, and since both values of $x^3$ are roots of unity, we can simplify the expression we want satisfied to the expression to $x^{1011}=x^3$. Since both values of $x^3$ are also values of $x^{1011}$, the roots for our $x^6+x^3+1$ are also roots of $x^{2022}+x^{1011}+1$, meaning that 
 \[x^6+x^3+1 | x^{2022}+x^{1011}+1\] so Therefore, the answer is $\boxed{E}$.

 - DavidHovey

 



\textbf{Solution 4 (Describe the Roots)}

We know that a monic polynomial $q$ divides a monic polynomial $p$ if and only if all the roots of $q$ are roots of $p.$ Since 
 \[P(x)=x^{2022}+x^{1011}+1=\frac{x^{3033}-1}{x^{1011}-1}\], the roots of $P$ are the $3033$rd roots of unity that aren't $1011$th roots of unity. 

 Now, note that:

 1: The roots of polynomial $A$ are the primitive $6$th roots of unity.

 2: The roots of polynomial $B$ are the primitive cube roots of unity.

 3: The roots of polynomial $C$ are the primitive $8$th roots of unity.

 4: The roots of polynomial $D$ are the primitive $18$th roots of unity.

 5: The roots of polynomial $E$ are the primitive $9$th roots of unity.

 However, since $6$, $8$, and $18$ don't divide $3033$, the roots of polynomial $A$ are not all $3033$rd roots of unity, and the same is true for polynomials $C$ and $D$, eliminating choices $A$, $C$ and $D.$ Also, since $3$ divides $1011$, the roots of polynomial $B$ are all $1011$th roots of unity, eliminating choice $B.$ That leaves choice $\boxed{E}$, and we can confirm that this is correct by noticing that $9$ divides $3033$ but not $1011.$ From that, we can see that the roots of polynomial $E$ are $3033$rd roots of unity but not $1011$th roots of unity, so they are all roots of $P.$ Therefore, $E$ divides $P.$




\textbf{Solution 5 (Simple Elimination)}

Put the value $x = -1$. This gives $P(-1) = (-1)^{2022} + (-1)^{1011} + 1 = 1$.This automatically eliminates choices $A$ , $C$ and $D$ since they do not form a factor of $1$ at $x = -1$.

 
Solution $5.1$:Now, put $x = \omega$ (omega) at choice $B$. We get that,
 \[\omega^2 + \omega + 1 = 0\]Thus choice B has $(x - \omega)$ as a factor. If choice $B$ were to be a factor of $P(x)$, $(x - \omega)$ would also have to be a factor of $P(x)$, which is clearly not the case, as $P(\omega) \neq 0$.This eliminates choice $B$, leaving us with answer $\boxed{E}$.

 
Solution $5.2$(Mods):Plug in $2$ for $x$ at choice $B$. $2^2+2*2+1=7$.The cycle of mods for $2$ and $7$ is $2$, $4$, $1$. Both $2022$ and $1011$ are divisible by $3$, so we have $1$ + $1$ + $1$ = $3$. $3$ is not divisible by $7$, so we eliminate choice $B$. This leaves us with $\boxed{(E) x^{6}+x^{3}+1}$.

 Solution $5.2$ ~ Bread 10

 



\textbf{Solution 5a (Elimination but slightly different)}

Like Solution 5, let $x=-1$ which eliminates the choices of $A$, $C$, and $D$ as they do not divide $P(-1)=1$ as they form $3, 2, 3$ respectively by letting $x=-1$.

 This leaves us with only $2$ choices, $B$ and $E$. Notice that letting $x=0$ or $x=1$ still make these answer choices work and the other values will leave large numbers for us to check which is not feasible in a 75 minutes math competition. 

 However, we notice answer choice $B$ is quadratic so if answer choice $B$ divides the given polynomial, then the roots of the quadratic must also be roots of the polynomial.

 Through quadratic formula, we find the roots of this quadratic as $\dfrac{-1+i\sqrt{3}}{2}$ and $\dfrac{-1-i\sqrt{3}}{2}$.

 We notice that these roots can be written nicely in polar form $\cos(\dfrac{2\pi}{3})+i\sin(\dfrac{2\pi}{3})$ or $\cos(\dfrac{4\pi}{3})+i\sin(\dfrac{4\pi}{3})$.

 We plug either one of these (preferably the root on the left as it is smaller) and see that the polynomial doesn't equal $0$ suggesting that $B$ is not the correct answer choice.

 As we only have one answer choice left, we choose $\boxed{E}$




\textbf{Solution 5b (elimination but even faster and cheesier)}

We first substitute $1$ and $-1$ and eliminate $A$, $C$, and $D$. We are now left with $B$ and $E$. When $x=3$, $B=13$; we then calculate $3^{2022}+3^{1011}+1$ mod $13$. Note that $2022$ and $1011$ are both multiples of $3$, and $3^3=27\equiv1\pmod{13}$. Thus, $3^{2022}\equiv3^{1011}\equiv3^3\equiv1\pmod{13}$, so it turns out that the given expression is $1+1+1=3$ mod $13$. We need it to be $0$, and since $3\neq0$, then $B$ is not the correct answer; $\boxed{\textbf{(E) }x^6+x^3+1}$ is correct. 

 


	\item (\textit{AMC 12 12A 2022}) Let $c$ be a real number, and let $z_1$ and $z_2$ be the two complex numbers satisfying the equation
$z^2 - cz + 10 = 0$. Points $z_1$, $z_2$, $\frac{1}{z_1}$, and $\frac{1}{z_2}$ are the vertices of (convex) quadrilateral $\mathcal{Q}$ in the complex plane. When the area of $\mathcal{Q}$ obtains its maximum possible value, $c$ is closest to which of the following?


 $\textbf{(A) }4.5 \qquad\textbf{(B) }5 \qquad\textbf{(C) }5.5 \qquad\textbf{(D) }6\qquad\textbf{(E) }6.5$

 \textbf{Answer: }A



\textbf{Solution 1}

Because $c$ is real, $z_2 = \bar z_1$.We have
 \begin{align*} 10 & = z_1 z_2 \\ & = z_1 \bar z_1 \\ & = |z_1|^2 , \end{align*}
where the first equality follows from Vieta's formula.

 Thus, $|z_1| = \sqrt{10}$.

 We have
 \begin{align*} c & = z_1 + z_2 \\ & = z_1 + \bar z_1 \\ & = 2 {\textrm Re}(z_1), \end{align*}
where the first equality follows from Vieta's formula.

 Thus, ${\textrm Re}(z_1) = \frac{c}{2}$.

 We have
 \begin{align*} \frac{1}{z_1} & = \frac{1}{10}\cdot\frac{10}{z_1} \\ & = \frac{1}{10}\cdot\frac{z_1 z_2}{z_1} \\ & = \frac{z_2}{10} \\ & = \frac{\bar z_1}{10}. \end{align*}
where the second equality follows from Vieta's formula.

 We have
 \begin{align*} \frac{1}{z_2} & = \frac{1}{10}\cdot\frac{10}{z_2} \\ & = \frac{1}{10}\cdot\frac{z_1 z_2}{z_2} \\ & = \frac{z_1}{10}. \end{align*}
where the second equality follows from Vieta's formula.

 Therefore,
 \begin{align*} {\textrm Area} \ Q & = \frac{1}{2} \left| {\textrm Re}(z_1) \right| \cdot 2 \left| {\textrm Im}(z_1) \right| \cdot \left( 1 - \frac{1}{10^2} \right) \\ & = \frac{1}{2} |c| \sqrt{10 - \frac{c^2}{4}} \left( 1 - \frac{1}{10^2} \right) \\ & = \frac{1 - \frac{1}{10^2}}{4} \sqrt{c^2 \left( 40 - c^2 \right)} \\ & \leq \frac{1 - \frac{1}{10^2}}{4} \cdot \frac{c^2 + \left( 40 - c^2 \right)}{2} \\ & = \frac{1 - \frac{1}{10^2}}{4} \cdot 20 , \end{align*}
where the inequality follows from the AM-GM inequality, and it is augmented to an equality if and only if $c^2 = 40 - c^2$.Thus, $|c| = 2 \sqrt{5} \approx \boxed{\textbf{(A) 4.5}}$.

 



\textbf{Solution 2 (Trapezoid)}

Since $c$, which is the sum of roots $z_1$ and $z_2$, is real, $z_1=\overline{z_2}$. 

 Let $z_1=a+bi$. Then $z_2=a-bi$. Note that the product of the roots is $10$ by Vieta's, so $z_1z_2=(a+bi)(a-bi)=a^2+b^2=10$.

 Thus, $\frac{1}{z_2}=\frac{1}{a-bi}\cdot\frac{a+bi}{a+bi}=\frac{a+bi}{a^2+b^2}=\frac{a+bi}{10}$. With the same process, $\frac{1}{z_1}=\frac{a-bi}{10}$.

 So, our four points are $a+bi,\frac{a+bi}{10},a-bi,$ and $\frac{a-bi}{10}$. WLOG let $a+bi$ be in the first quadrant and graph these four points on the complex plane. Notice how quadrilateral Q is a trapezoid with the real axis as its axis of symmetry. It has a short base with endpoints $\frac{a+bi}{10}$ and $\frac{a-bi}{10}$, so its length is $\left|\frac{a+bi}{10} - \frac{a-bi}{10}\right| =\frac{b}{5}$. Likewise, its long base has endpoints $a+bi$ and $a-bi$, so its length is $|(a+bi)-(a-bi)|=2b$. 

 The height, which is the distance between the two lines, is the difference between the real values of the two bases $\implies h= a-\frac{a}{10}=\frac{9a}{10}$.

 Plugging these into the area formula for a trapezoid, we are trying to maximize $\frac12\cdot\left(2b+\frac{b}{5}\right)\cdot\frac{9a}{10}=\frac{99ab}{100}$. Thus, the only thing we need to maximize is $ab$.

 With the restriction that $a^2+b^2=(a-b)^2+2ab=10$, $ab$ is maximized when $a=b=\sqrt{5}$. 

 Remember, $c$ is the sum of the roots, so $c=z_1+z_2=2a=2\sqrt5=\sqrt{20}\approx\boxed{4.5}$




\textbf{Solution 3 (Fast)}

Like the solutions above we can know that $|z_1| = |z_2| = \sqrt{10}$ and $z_1=\overline{z_2}$.

 Let $z_1=\sqrt{10}e^{i\theta}$ where $0<\theta<\pi$, then $z_2=\sqrt{10}e^{-i\theta}$, $\frac{1}{z_1}=\frac{1}{\sqrt{10}}e^{-i\theta}$, $\frac{1}{z_2}= \frac{1}{\sqrt{10}}e^{i\theta}$.

 On the basis of symmetry, the area $A$ of $\mathcal{Q}$ is the difference between two isoceles triangles,so

 $2A=10\sin2\theta-\frac{1}{10}\sin2\theta\leq10-\frac{1}{10}$. The inequality holds when $2\theta=\frac{\pi}{2}$, or $\theta=\frac{\pi}{4}$.

 Thus, $c= 2 {\textrm Re} \ z_1 =2 \sqrt{10} \cos\frac{\pi}{4}=\sqrt{20} \approx \boxed{\textbf{(A) 4.5}}$.

 



\textbf{Solution 4 (Calculus Finish)}

Like in Solution 3, we find that $Q = \frac12\cdot\left(2b+\frac{b}{5}\right)\cdot\frac{9a}{10}= \frac{99}{100}ab$, thus, $Q$ is maximized when $ab$ is maximized. $ab = a \cdot \sqrt{10 - a^2} = \sqrt{10a^2 - a^4}$, let $f(a) =  \sqrt{10a^2 - a^4}$.

 By the Chain Rule and the Power Rule, $f'(a) = \frac12 \cdot (10a^2 - a^4)^{-\frac12} \cdot (10(2a)-4a^3)) = \frac{20a-4a^3}{ 2\sqrt{10a^2 - a^4} } = \frac{10a-2a^3}{ \sqrt{10a^2 - a^4} }$
$f'(a) = 0$, $10a-2a^3 = 0$, $a \neq 0$, $a^2 = 5$, $a = \sqrt{5}$
$\because f'(a) = 0$ when $a = \sqrt{5}$, $f'(a)$ is positive when $a < \sqrt{5}$, and $f'(a)$ is negative when $a > \sqrt{5}$
$\therefore f(a)$ has a local maximum when $a = \sqrt{5}$.

 Notice that $ab = \frac{ (a+b)^2 - (a^2+b^2) }{2} = \frac{c^2 - 10}{2}$, $\frac{c^2 - 10}{2} = 5$, $c = \sqrt{2 \cdot 5 + 10} = \sqrt{20} \approx \boxed{\text{(A) 4.5}}$




\textbf{Solution 5 (calculus but it's bash)}

Note that $z=\dfrac c2\pm\dfrac{\sqrt{c^2-40}}2$, so let $z_1=\dfrac c2+\dfrac{\sqrt{c^2-40}}2$ and $z_2=\dfrac c2-\dfrac{\sqrt{c^2-40}}2$. Taking a look at the answer choices, they range between $c=4.5$ to $c=6.5$, and in that range, $c^2$ is always less than $40$. Thus, $c^2-40<0$ for our possible answer choices; we can then rewrite $z_1$ and $z_2$ as $\dfrac c2+\dfrac{\sqrt{40-c^2}}2i$ and $\dfrac c2-\dfrac{\sqrt{40-c^2}}2i$, respectively, with real coefficients. 

 Let us compute $\dfrac1z$: 

 
 \[\dfrac1z=\dfrac1{\frac c2\pm\frac{\sqrt{c^2-40}}2}=\dfrac{\frac c2\mp\frac{\sqrt{c^2-40}}2}{\left(\frac c2\right)^2-\left(\frac{\sqrt{c^2-40}}2\right)^2}=\dfrac{\frac c2\mp\frac{\sqrt{c^2-40}}2}{\frac{c^2}4-\frac{c^2-40}4}=\dfrac{2c\mp2\sqrt{c^2-40}}{40}=\dfrac{c\mp\sqrt{c^2-40}}{20}.\] 

 Then, $\dfrac1{z_1}=\dfrac{c-\sqrt{c^2-40}}{20}=\dfrac c{20}-\dfrac{\sqrt{40-c^2}}{20}i$ while $\dfrac1{z_2}=\dfrac{c+\sqrt{c^2-40}}{20}=\dfrac c{20}+\dfrac{\sqrt{40-c^2}}{20}i$. 

 In the complex plane, we can draw a rough sketch of $z_1,z_2,\dfrac1{z_1},\dfrac1{z_2}$: 

 
                    \begin{figure}[H]
                        \centering
                        \includegraphics[width=0.5\textwidth]{Solution/AMC_12_12A_2022_22_Solution_Figure_1.png}
                    \end{figure}
                    

 Note that, here, our quadrilateral is an isoceles trapezoid. The shorter base length is $\left(\dfrac{\sqrt{40-c^2}}{20}-\left(-\dfrac{\sqrt{40-c^2}}{20}\right)\right)=\dfrac1{10}\sqrt{40-c^2}$. 

 The longer base length is $\left(\dfrac{\sqrt{40-c^2}}2-\left(-\dfrac{\sqrt{40-c^2}}2\right)\right)=\sqrt{40-c^2}$. 

 The average of the two bases is $\dfrac{11}{20}\sqrt{40-c^2}$. 

 The height of our trapezoid (which is horizontal parallel to the $x$-axis in our diagram above) is simply $\dfrac c2-\dfrac c{20}=\dfrac9{20}c$. 

 Since the area of a trapezoid is the product of the average of its bases and its height, we conclude that our trapezoid's area is $\dfrac{99}{400}c\sqrt{40-c^2}$, which is a function of $c$. Thus, let $A(c)=\dfrac{99}{400}c\sqrt{40-c^2}$. 

 Taking the derivative (FINALLY we get to the Calculus part haha, this definitely wasn't clickbait :3 trust me) with respect to $c$, we find that $\dfrac{dA}{dc}=\dfrac{99}{400}\left(\sqrt{40-c^2}+c\left(\dfrac{-2c}{2\sqrt{40-c^2}}\right)\right)=\dfrac{99}{400}\left(\sqrt{40-c^2}-\dfrac{c^2}{\sqrt{40-c^2}}\right)$. 

 To find an extremum, we set the derivative equal to zero: 

 \begin{align*}\dfrac{dA}{dc}&=\dfrac{99}{400}\left(\sqrt{40-c^2}-\dfrac{c^2}{\sqrt{40-c^2}}\right) \\0&=\dfrac{99}{400}\left(\sqrt{40-c^2}-\dfrac{c^2}{\sqrt{40-c^2}}\right) \\\sqrt{40-c^2}-\dfrac{c^2}{\sqrt{40-c^2}}&=0 \\\sqrt{40-c^2}&=\dfrac{c^2}{\sqrt{40-c^2}} \\\left(\sqrt{40-c^2}\right)^2&=c^2 \\c^2&=40-c^2 \\2c^2&=40 \\c^2&=20 \\c&=\sqrt{20} \\&\approx4.47 \\\end{align*}


 Clearly, this is very close to $\boxed{\textbf{(A)}~4.5}$, so we are done. QED. 

 



\textbf{Solution 6 (AM-GM Inequality)}

First and foremost, because $z_n$ is a complex number, and the value of $z$ could be found in terms of $c$, the given quadratic equation could be solved. WLOG, $z_1=\frac{c+\sqrt{c^2-40}}{2}$ and $z_1=\frac{c-\sqrt{c^2-40}}{2}$. If $c^2-40\ge0$, no quadrilateral will form. Thereby, it could be inferred that $c^2-40<0$.

 Each complex number could be represented as a point.
 \begin{align*}     &z_1\left(\frac{c}{2}, \frac{\sqrt{40-c^2}}{2}\right) \\     &z_2\left(\frac{c}{2}, \frac{-\sqrt{40-c^2}}{2}\right) \\     &\frac{1}{z_1}\left(\frac{c}{20}, \frac{c-\sqrt{40-c^2}}{20}\right) \\     &\frac{1}{z_2}\left(\frac{c}{20}, \frac{c+\sqrt{40-c^2}}{20}\right) \end{align*}
Notice that vertices form a trapezoid. Therefore, the area of the trapezoid $Q$ could be represented as
 \[\frac{\left(\frac{c}{2}-\frac{c}{20}\right)\left\{\left(\frac{\sqrt{40-c^2}}{2}+\frac{\sqrt{40-c^2}}{2}\right)+\left(\frac{c+\sqrt{40-c^2}}{20}-\frac{c-\sqrt{40-c^2}}{20}\right)\right\}}{2}.\]Since we are finding the maximum value, any desired constant could be multiplied. Multiplying 80 and other constants to the equation above may provide a simplified form.
 \begin{align*} &(9c)(20\sqrt{40-c^2}+2\sqrt{40-c^2}) \\ \Rightarrow&9c(22\sqrt{40-c^2}) \\ \Rightarrow&c\sqrt{40-c^2} \\ \end{align*}
We can use AM-GM Inequality! According to the AM-GM Inequality,
 \[40-c^2+c^2\ge2\sqrt{(40-c^2)\cdot c^2}.\]$\sqrt{(40-c^2)\cdot c^2}$ reaches maximum when $40-c^2=c^2$. In another words, $c=\pm2\sqrt{5}$. Using $2.2$ for the approximation of $\sqrt{5}$, $\max(c)\approx4.4$. Thus, the closest value in the answer choice is $\boxed{\textbf{(A) }4.5}$.

 


	\item (\textit{AMC 12 12A 2022}) Let $h_n$ and $k_n$ be the unique relatively prime positive integers such that 
 \[\frac{1}{1}+\frac{1}{2}+\frac{1}{3}+\cdots+\frac{1}{n}=\frac{h_n}{k_n}.\] Let $L_n$ denote the least common multiple of the numbers $1, 2, 3, \ldots, n$. For how many integers with $1\le{n}\le{22}$ is $k_n<L_n$?


 $\textbf{(A) }0 \qquad\textbf{(B) }3 \qquad\textbf{(C) }7 \qquad\textbf{(D) }8\qquad\textbf{(E) }10$

 \textbf{Answer: }D



\textbf{Solution 1}

We are given that 
 \[\sum_{i=1}^{n}\frac1i = \frac{1}{L_n}\sum_{i=1}^{n}\frac{L_n}{i} = \frac{h_n}{k_n}.\] Since $k_n < L_n,$ we need $\gcd\left(\sum_{i=1}^{n}\frac{L_n}{i}, L_n\right)>1.$

 For all primes $p$ such that $p\leq n,$ let $\nu_p(L_n)=e\geq1$ be the exponent of the largest power of $p$ that divides $L_n.$

 It is clear that $L_n\equiv0\pmod{p},$ so we test whether $\sum_{i=1}^{n}\frac{L_n}{i}\equiv0\pmod{p}.$ Note that 
 \[\sum_{i=1}^{n}\frac{L_n}{i} \equiv \sum_{i=1}^{\left\lfloor\frac{n}{p^e}\right\rfloor}\frac{L_n}{ip^e} \ (\operatorname{mod} \ p^e) \equiv \sum_{i=1}^{\left\lfloor\frac{n}{p^e}\right\rfloor}\frac{L_n}{ip^e} \ (\operatorname{mod} \ p).\]We construct the following table for $\nu_p(L_n)=e:$
 \[\begin{array}{c|c|l|c}  \textbf{Case of }\boldsymbol{(p,e)} & \textbf{Interval of }\boldsymbol{n} & \hspace{22.75mm}\textbf{Sum} & \boldsymbol{\stackrel{?}{\equiv}0\ (\operatorname{mod} \ p)} \\ [0.5ex] \hline\hline  & & & \\ [-2ex] (2,1) & [2,3] & L_n/2 &  \\      (2,2) & [4,7] & L_n/4 &  \\     (2,3) & [8,15] & L_n/8 & \\     (2,4) & [16,22] & L_n/16 &  \\ [0.5ex] \hline   & & & \\ [-2ex] (3,1) & [3,5] & L_n/3 & \\     & [6,8] & L_n/3 + L_n/6 & \checkmark \\    (3,2) & [9,17] & L_n/9 & \\  & [18,22] & L_n/9 + L_n/18 & \checkmark \\ [0.5ex] \hline   & & & \\ [-2ex] (5,1) & [5,9] & L_n/5 & \\     & [10,14] & L_n/5 + L_n/10 & \\ & [15,19] & L_n/5 + L_n/10 + L_n/15 & \\   & [20,22] & L_n/5 + L_n/10 + L_n/15 + L_n/20 & \checkmark \\ [0.5ex] \hline   & & & \\ [-2ex] (7,1) & [7,13] & L_n/7 & \\     & [14,20] & L_n/7 + L_n/14 & \\ & [21,22] & L_n/7 + L_n/14 + L_n/21 & \\ [0.5ex] \hline   & & & \\ [-2ex] (11,1) & [11,21] & L_n/11 &  \\ & \{22\} & L_n/11 + L_n/22 & \\ [0.5ex] \hline & & & \\ [-2ex] (13,1) & [13,22] & L_n/13 &  \\ [0.5ex] \hline & & & \\ [-2ex] (17,1) & [17,22] & L_n/17 &  \\ [0.5ex] \hline & & & \\ [-2ex] (19,1) & [19,22] & L_n/19 & \\ [0.5ex] \end{array}\]Note that:

 
            \begin{enumerate} 
        \item If the \textbf{Sum} column has only one term, then it is never congruent to $0$ modulo $p.$

        \item If $p$ and $q$ are positive integers such that $p\geq q,$ then $L_p$ is a multiple of $L_q.$ Therefore, for a specific case, if the sum is congruent to $0$ modulo $p$ for the smallest element in the interval of $n,$ then it is also congruent to $0$ modulo $p$ for all other elements in the interval of $n.$

            \end{enumerate}
            Together, there are $\boxed{\textbf{(D) }8}$ such integers $n,$ namely 
 \[6,7,8,18,19,20,21,22.\]~MRENTHUSIASM

 



\textbf{Solution 2}

Notice that $\sum_{i=1}^{n}\frac{1}{i} = \frac{\sum_{i=1}^{n}\frac{L_n}{i}}{L_n}$. Thus, in order for $L_n$ to reduce to a smaller value, the numerator and denominator must share a common factor. 

 We can start by testing some values of $n$, and quickly observe that $n=6$ is the first number that satisfies the desired conditions. In particular, $\sum_{i=1}^{6}\frac{1}{i} = \frac{\sum_{i=1}^{6}\frac{L_6}{i}}{L_6} = \frac{147}{60} = \frac{49}{20}$. Since a factor of three is shared, we are motivated to observe factors of three in the numerator $\sum_{i=1}^{6}\frac{60}{i}$. Notice that $\nu_3(L_6) = 1$, since $3^1$ is the greatest power of three that occurs in $1,2,3,4,5,6$. Consider the sum in the numerator in mod 3. We have $(\frac{60}{1}, \frac{60}{2}, \frac{60}{3}, \frac{60}{4}, \frac{60}{5}, \frac{60}{6} )\rightarrow (0,0,2,0,0,1)$. Adding 0+0+2+0+0+1 we get 0 mod 3, which is why we are able to simplify $\frac{147}{60}$.

 Thus, we should be taking the numerator mod p, where p is a prime. 

 Quickly, we wish to show that the numerator and denominator can never share a factor of 2. Say we have least common multiple $L_n$. Let $\nu_2(L_n) = e$. In other words, let $2^e$ be the largest power of 2 such that $2^e$ divides $L_n$. Notice that among all numbers $1$ through $n$, only one of them can be a multiple of $2^e$, being $2^e$ itself. The next multiple of $2^e$ would be $2 \cdot 2^e = 2^{e+1}$, contradicting the fact that $\nu_2(L_n) = e$. This means that in the sum $\sum_{i=1}^{n}\frac{L_n}{i}$, there will be $1$ value that is $1$ mod $2$, resulting from $\frac{L_n}{2^e}$, and the rest are all $0$s. Clearly, the numerator will then be $1$ mod $2$, which is not a multiple of 2.

 When $p=3$, $n=6,7,8,18,19,20,21,22$ will all work, and this can be easily tested through taking the numerator mod $3$. $p=5$ is only satisfied when the $L_n$ contains $(5,10,15,20)$, so when $n\ge 20$. However, this is just an overlap with the previous values, so there are $\boxed{\textbf{(D) }8}$.

 


	\item (\textit{AMC 12 12A 2022}) How many strings of length $5$ formed from the digits $0$, $1$, $2$, $3$, $4$ are there such that for each $j \in \{1,2,3,4\}$, at least $j$ of the digits are less than $j$? (For example, $02214$ satisfies this condition
because it contains at least $1$ digit less than $1$, at least $2$ digits less than $2$, at least $3$ digits less
than $3$, and at least $4$ digits less than $4$. The string $23404$ does not satisfy the condition because it
does not contain at least $2$ digits less than $2$.)


 $\textbf{(A) }500\qquad\textbf{(B) }625\qquad\textbf{(C) }1089\qquad\textbf{(D) }1199\qquad\textbf{(E) }1296$

 \textbf{Answer: }E



\textbf{Solution 1 (Parking Functions)}

For some $n$, let there be $n+1$ parking spaces counterclockwise in a circle. Consider a string of $n$ integers $c_1c_2 \ldots c_n$ each between $0$ and $n$, and let $n$ cars come into this circle so that the $i$th car tries to park at spot $c_i$, but if it is already taken then it instead keeps going counterclockwise and takes the next available spot. After this process, exactly one spot will remain empty.

 Then the strings of $n$ numbers between $0$ and $n-1$ that contain at least $k$ integers $<k$ for $1 \leq k \leq n$ are exactly the set of strings that leave spot $n$ empty. Also note for any string $c_1c_2 \ldots c_n$, we can add $1$ to each $c_i$ (mod $n+1$) to shift the empty spot counterclockwise, meaning for each string there exists exactly one $j$ with $0 \leq j \leq n$ so that $(c_1+j)(c_2+j) \ldots (c_n+j)$ leaves spot $n$ empty. This gives there are $\frac{(n+1)^{n}}{n+1} = (n+1)^{n-1}$ such strings. 

 Plugging in $n = 5$ gives $\boxed{\textbf{(E) }1296}$ such strings.

 



\textbf{Solution 2 (OmegaLearn)}

Consider some string that satisfies the problem statement, with $a$ 0's, $b$ 1's, $c$ 2's, $d$ 3's, and $e$ 4's, so $a+b+c+d+e=5$. There are exactly $\dfrac{5!}{a!b!c!d!e!}$ ways to order the given string. This number will be divisible by 5 unless one of $a, b, c, d, e$ is $5$, which implies that the rest are $0$. This corresponds to the string having $5$ of one digit. But the only possible string that has $1$ distinct digit is $00000$. Thus, the number of strings must be $1 \pmod{5}$. The only answer choice that satisfies this is $\boxed{\textbf{(E) }1296}$.

 



\textbf{Solution 3 (Casework)}

Note that a valid string must have at least one $0.$

 We perform casework on the number of different digits such strings can have. For each string, we list the digits in ascending order, then consider permutations:

 
            \begin{enumerate} 
        \item The string has $1$ different digit.The only possibility is $00000.$ 

 \textbf{There is $\boldsymbol{1}$ string in this case.} 

 
        \item The string has $2$ different digits.We have the following table:
 \[\begin{array}{c||c|c|c|c||c} & & & & & \\ [-2.5ex] \textbf{Digits} & \boldsymbol{01} & \boldsymbol{02} & \boldsymbol{03} & \boldsymbol{04} & \textbf{Row's Count} \\ [0.5ex] \hline & & & & & \\ [-1.5ex] & 00001 & 00002 & 00003 & 00004 & \hspace{2mm}4\cdot\frac{5!}{4!1!}=20 \\ [2ex]  & 00011 & 00022 & 00033 & & \hspace{2mm}3\cdot\frac{5!}{3!2!}=30 \\ [2ex]  & 00111 & 00222 & & & \hspace{2mm}2\cdot\frac{5!}{2!3!}=20 \\ [2ex]  & 01111 & & & & 1\cdot\frac{5!}{1!4!}=5 \\ [0.75ex] \end{array}\]\textbf{There are $\boldsymbol{20+30+20+5=75}$ strings in this case.} 

 
        \item The string has $3$ different digits.We have the following table:
 \[\begin{array}{c||c|c|c|c|c|c||c} & & & & & & &  \\ [-2.5ex] \textbf{Digits} & \boldsymbol{012} & \boldsymbol{013} & \boldsymbol{014} & \boldsymbol{023} & \boldsymbol{024} & \boldsymbol{034} & \textbf{Row's Count} \\ [0.5ex] \hline & & & & & & & \\ [-1.5ex] & 00012 & 00013 & 00014 & 00023 & 00024 & 00034 & \hspace{2mm}6\cdot\frac{5!}{3!1!1!}=120 \\ [2ex]  & 00112 & 00113 & 00114 & 00223 & 00224 & & \hspace{2mm}5\cdot\frac{5!}{2!2!1!}=150 \\ [2ex]  & 00122 & 00133 & & 00233 & & & 3\cdot\frac{5!}{2!1!2!}=90 \\ [2ex]  & 01112 & 01113 & 01114 & & & & 3\cdot\frac{5!}{1!3!1!}=60 \\ [2ex] & 01122 & 01133 & & & & & 2\cdot\frac{5!}{1!2!2!}=60 \\ [2ex] & 01222 & & & & & & 1\cdot\frac{5!}{1!1!3!}=20 \\ [0.75ex] \end{array}\]\textbf{There are $\boldsymbol{120+150+90+60+60+20=500}$ strings in this case.} 

 
        \item The string has $4$ different digits.We have the following table:
 \[\begin{array}{c||c|c|c|c} & & & & \\ [-2.5ex] \textbf{Digits} & \boldsymbol{0123} & \boldsymbol{0124} & \boldsymbol{0134} & \boldsymbol{0234} \\ [0.5ex] \hline & & & & \\ [-1.5ex] & 00123 & 00124 & 00134 & 00234 \\ [2ex]  & 01123 & 01124 & 01134 & \\ [2ex]  & 01223 & 01224 & & \\ [2ex]  & 01233 & & & \\ [0.75ex] \end{array}\]\textbf{There are $\boldsymbol{10\cdot\frac{5!}{2!1!1!1!}=600}$ strings in this case.} 

 
        \item The string has $5$ different digits.\textbf{There are $\boldsymbol{5!=120}$ strings in this case.}

 
            \end{enumerate}
            Together, the answer is $1+75+500+600+120=\boxed{\textbf{(E) }1296}.$




\textbf{Solution 4 (Recursive Equations Approach)}

Denote by $N \left( p, q \right)$ the number of $p$-digit strings formed by using numbers $0, 1, \cdots, q$, where for each $j \in \{1,2, \cdots , q\}$, at least $j$ of the digits are less than $j$.

 We have the following recursive equation:
 \[N \left( p, q \right) = \sum_{i = 0}^{p - q} \binom{p}{i} N \left( p - i, q - 1 \right) , \ \forall \ p \geq q \mbox{ and } q \geq 1\]and the boundary condition $N \left( p, 0 \right) = 1$ for any $p \geq 0$.

 By solving this recursive equation, for $q = 1$ and $p \geq q$, we get
 \begin{align*} N \left( p , 1 \right) & = \sum_{i = 0}^{p - 1} \binom{p}{i} N \left( p - i, 0 \right) \\ & = \sum_{i = 0}^{p - 1} \binom{p}{i} \\ & = \sum_{i = 0}^p \binom{p}{i} - \binom{p}{p} \\ & = 2^p - 1 . \end{align*}


 For $q = 2$ and $p \geq q$, we get
 \begin{align*} N \left( p , 2 \right) & = \sum_{i = 0}^{p - 2} \binom{p}{i} N \left( p - i, 1 \right) \\ & = \sum_{i = 0}^{p - 2} \binom{p}{i} \left( 2^{p - i} - 1 \right) \\ & = \sum_{i = 0}^p \binom{p}{i} \left( 2^{p - i} - 1 \right) - \sum_{i = p - 1}^p \binom{p}{i} \left( 2^{p - i} - 1 \right) \\ & = \sum_{i = 0}^p \left( \binom{p}{i} 1^i 2^{p - i} - \binom{p}{i} 1^i 1^{p - i} \right) - p \\ & = \left( 1 + 2 \right)^p - \left( 1 + 1 \right)^p - p \\ & = 3^p - 2^p - p . \end{align*}


 For $q = 3$ and $p \geq q$, we get
 \begin{align*} N \left( p , 3 \right) & = \sum_{i = 0}^{p - 3} \binom{p}{i} N \left( p - i, 2 \right) \\ & = \sum_{i = 0}^{p - 3} \binom{p}{i} \left( 3^{p - i} - 2^{p - i} - \left( p - i \right) \right) \\ & = \sum_{i = 0}^p \binom{p}{i} \left( 3^{p - i} - 2^{p - i} - \left( p - i \right) \right) - \sum_{i = p - 2}^p \binom{p}{i} \left( 3^{p - i} - 2^{p - i} - \left( p - i \right) \right) \\ & = \sum_{i = 0}^p \left( \binom{p}{i} 1^i 3^{p - i} - \binom{p}{i} 1^i 2^{p - i} - \binom{p}{i} \left( p - i \right) \right) - \frac{3}{2} p \left( p - 1 \right) \\ & = \left( 1 + 3 \right)^p - \left( 1 + 2 \right)^p - \frac{d \left( 1 + x \right)^p}{dx} \bigg|_{x = 1} - \frac{3}{2} p \left( p - 1 \right) \\ & = 4^p - 3^p - 2^{p-1} p - \frac{3}{2} p \left( p - 1 \right) . \end{align*}


 For $q = 4$ and $p = 5$, we get
 \begin{align*} N \left( 5 , 4 \right) & = \sum_{i = 0}^1 \binom{5}{i} N \left( 5 - i , 3 \right) \\ & = N \left( 5 , 3 \right) + 5 N \left( 4 , 3 \right) \\ & = \boxed{\textbf{(E) }1296}  . \end{align*}





\textbf{Solution 6 (Observations)}

The number of strings is $(n+1)^{(n-1)}$ as shown by Solution 1 (Parking Function), which is always equivalent to 1 (mod n). Thus you can choose $\boxed{\textbf{(E) }1296}$.

 However, you **CANNOT** prove that fact from this *incorrect* argument:

 Let the set of all valid sequences be $S$. Notice that for any sequence $\{a_1,a_2,a_3,a_4,a_5\}$ in $S$, the 5 cyclic permutations
 \begin{align*} \{a_1,a_2,a_3,a_4,a_5\}\\ \{a_2,a_3,a_4,a_5,a_1\}\\ \{a_3,a_4,a_5,a_1,a_2\}\\ \{a_4,a_5,a_1,a_2,a_3\}\\ \{a_5,a_1,a_2,a_3,a_4\} \end{align*}
must also belong in $S$. However, one must consider the edge case all 5 elements are the same (only $\{0,0,0,0,0\}$), in which case all sequences listed are equivalent.

 Alas, one must also consider cases like 10101 that only have 3 distinct cyclic permutations, so we cannot get a useful constraint from this view, unless *all* the cases are categorized and counted. 

 Thus it is NOT implied that $\lvert S \rvert \equiv 1  \pmod 5$, so it is not justified to choose $\boxed{\textbf{(E) }1296}$ by inspection of answer options.

 If you solve $S(n)$ for smaller $n$ by bashing, you can find $|S(1)| =1$, $|S(2)| = 3$, $|S(3)| =16$, $|S(4)| =125$, which leads to an Engineer's Induction guess that $|S(n)| \equiv 1 \pmod n$, or the exact formula $(n+1)^{(n-1)}$

 Note: I'm pretty sure this solution works. Firstly, 10101 does have 5 distinct cyclic permutations, namely 10101, 11010, 01101, 10110, and 01011. Secondly, any sequence of 5 elements must have either 1 or a multiple of 5 distinct cyclic permutations. If a sequence of 5 elements has only 1 cyclic permutation, that must be the case with $(0, 0, 0, 0, 0)$, as noted (since 0 is necessary in our sequence when $j = 1$). If the sequence is cyclic every $k$ "rotations/shifts," then it is ONLY going to be at the starting sequence when you shift it $0, k, 2k, \dots$ times. Since shifting any sequence 5 times takes it back to itself, we must have $k \mid 5 \implies k = 5$ (since we already considered $k = 1$ with $(0, 0, 0, 0, 0)$), so the cycle must be $5$. So really, this solution does work, is not by coincidence, and is a very nice (albeit cheap) solution.

 TL;DR: The sentence "Alas, one must also consider cases like 10101 that only have 3 distinct cyclic permutations" is incorrect, for these cases do not actually exist. It's trivial to check 10101 has 5 cyclic permutations.

 



\textbf{Solution 5 (Mod 5 Trick Redeemed)}

Solution 4 tried to observe the answer modulo $5$ to easily solve the problem, but apparently had faulty logic. This solution is still completely viable though:

 Notice that for any valid set $\{a_1, a_2, a_3, a_4, a_5\}$, if there is at least one element in the set that is unique (i.e. there is at least one digit in the set that is found nowhere else in the set), then the number of distinct permutations of the set is clearly divisible by $5$ (alternatively there is no factor of 5 in the denominator of the multinomial coefficient unless an element is repeated 5 times). Therefore to evaluate the answer modulo $5$, we only need to look at sets where each element has a multiplicity of at least $2$ (i.e. appears twice or more in the set). 

 These sets are of the form $\{a,a,b,b,b\}$ and $\{a,a,a,a,a\}$. The first set can be permuted in $\binom{5}{2,3}=10 \equiv 0\pmod{5}$, and the second set can be permuted one way, and the only set of the form $\{a,a,a,a,a\}$ is $\{0,0,0,0,0\}$. Therefore the answer is congruent to $1\pmod{5}$ and you CAN choose $\boxed{\textbf{(E) }1296}$.

 



\textbf{Solution 7 (Casework Bash)}

We list out all possibilities for the string, sorted in increasing order. Note that 1 cannot appear before the 2nd position, 2 cannot appear before the 3rd position, and so on. We also make a list of the number of possible permutations of each sorted string.

 00000 - 1

 00001 - 5

 00002 - 5

 00003 - 5

 00004 - 5

 00011 - 10

 00012 - 20

 00013 - 20

 00014 - 20

 00022 - 10

 00023 - 20

 00024 - 20

 00033 - 10

 00034 - 20

 00111 - 10

 00112 - 30

 00113 - 30

 00114 - 30

 00122 - 30

 00123 - 60

 00124 - 60

 00133 - 30

 00134 - 60

 00222 - 10

 00223 - 30

 00224 - 30

 00233 - 30

 00234 - 60

 01111 - 5

 01112 - 20

 01113 - 20

 01114 - 20

 01122 - 30

 01123 - 60

 01124 - 60

 01133 - 30

 01134 - 60

 01222 - 20

 01223 - 60

 01224 - 60

 01233 - 60

 01234 - 120

 
Finally, adding all of these up, we obtain $1\cdot 1+5\cdot 5+5\cdot 10+10\cdot 20+10\cdot 30 + 10\cdot 60 + 1\cdot 120 = \boxed{\textbf{(E) }1296}$. (Or alternatively, we could notice as in solution 5 that the only possible string with only 1 permutation is 00000, and thus the answer must be congruent to 1 modulo 5)

 The key to organizing all of these cases is to increase the digits one by one, starting from the right and iterating through all the possible cases. 

 



\textbf{Solution 8 (Complementary Counting)}

Let us call the condition where you require at least $j$ digits less than $j$ Rule $j$.For a sequence to break Rule $1$, it must contain no $0$s. There are $4^5=1024$ ways to do this.For a sequence to break Rule $2$ but not Rule $1$, it must contain exactly 1 $0$. However, this one $0$ contributes toward completing the rule, meaning all other digits must be greater than $2$. After placing the one $0$ in any of $5$ slots, there are $4$ slots left with $3$ options each, giving a total of $5 \cdot 3^4 = 405$ ways.For a sequence to break Rule $3$ but neither Rule $2$ or $1$, it must have two digits less than $2$ (either $(0, 0)$ or $(0, 1)$) and all other digits must be $3$ or $4$. There are a total of $30$ ways to arrange the first two digits and $2^3$ ways to arrange the latter three, leaving $240$ ways to do this.Finally, for a sequence to only break rule $4$, it must have three digits less than $3$. They could be $(0, 1, 2)$ (60 arr.), $(0, 1, 1); (0, 0, 1); (0, 0, 2)$ ($30$ arr. each), or $(0, 0, 0)$ ($10$ arr.) This gives $160$ possible arrangements overall.

 There are in total $5^5=3125$ possible arrangements, so by adding all of the invalid states together and subtracting from $3125$, we get our answer, which is $\boxed{\textbf{(E) }1296}$

 Great idea! ~Bigbrain\_2009

 



\textbf{Solution 9 (Recursion)}

We will use recursion to solve this problem. Let us break up the problem's condition into criteria, with criterion \( k \) being the requirement that at least \( k \) of the digits are less than \( k \). Define \( A_n \) as the number of strings of length \( n \) formed from the digits \( 0, 1, 2, \dots, n-1 \) such that at least \( j \) of the digits are less than \( j \) for all integers \( j \); \( 1 \le j \le n-1 \).

 Note that the final criterion any string counted as apart of \( A_{n-1} \) is required to meet is criterion \( n-2 \). However, all strings that \( A_{n-1} \) counts will also meet criterion \( n-1 \). This is due to the maximum value of such strings' \( n-1 \) digits being \( n-2 \). Initially, we have \( A_1 = 1 \). We check cases and find \( A_2 = 3 \). To generate recursion, consider the compliment of the problem's condition. Without the condition, we are counting strings of length \( n \) with \( n \) choices for each digit. There are of course \( n^n \) such strings. To count the number of strings that do not meet at least one criterion, we denote a string as \( k \)-bad if the first criterion the string violates is criterion \( k \). For example, the string $33040$ is $3$-bad because it meets criteria 1 and 2, but violates criterion 3. Similarly, the string $13223$ is $1$-bad because it does contain 1 digit less than 1. We can count the number of strings that are \( k \)-bad by first choosing and assigning \( k-1 \) of the string's \( n \) spots to a string counted as apart of \( A_{k-1} \). This can be done in \( \binom{n}{k-1} \cdot A_{k-1} \) ways. This \( (k-1) \) digit string will meet all criteria up to criterion \( k-1 \). To ensure criterion \( k \) is not met, we fill up the remaining \( n-(k-1) \) spots with a digit from \( k \) to \( n-1 \), inclusive. This can be done in \( (n-k)^{n+1-k} \) ways. We sum over all integers \( k \), \( 1 \le k \le n-1 \) to count all strings that do not meet at least one criterion.

 
 \[\sum_{k=1}^{n-1} \binom{n}{k-1} A_{k-1} (n-k)^{n+1-k}\]This is the same as
 \[\sum_{k=2}^{n} \binom{n}{k-2} A_{k-2} (n-k+1)^{n+2-k}=\sum_{k=2}^{n} \binom{n}{n+2-k} A_{k-2} (n-k+1)^{n+2-k}\]We substitute \( n+2-k \) for \( k \) as the expression takes on the same set of values. We have
 \[\sum_{k=2}^{n} \binom{n}{n+2-(n+2-k)} A_{(n+2-k)-2} (n-(n+2-k)+1)^{n+2-(n+2-k)} = \sum_{k=2}^{n} \binom{n}{k} A_{n-k} (k-1)^k\]Therefore,
 \[A_n = n^n - \sum_{k=2}^{n} \binom{n}{k} A_{n-k} (k-1)^k\]Setting \( A_0 = 1 \) and \( A_1 = 1 \), we use the recursion to derive \( A_2 = 3 \), \( A_3 = 16 \), \( A_4 = 125 \), \( A_5 = 1296 \). In our problem, we were counting 5 digits strings. Thus, the answer is $\boxed{\textbf{(E) }1296}$




\textbf{Solution 10 (Parity Cheese)}

We note that there are 6 different multisets of digits that have an odd number of permutations, namely, 00000, 00001, 00002, 00003, 0004, 01111. We know that there cannot be any other possibilities since, if any digit appears 2 or 3 times, there are $\binom52=\binom53=10$ ways to place those two digits, and thus the total number of permutations must be even. Additionally, if there are at least two digits that appear once, there are $5\cdot 4=20$ ways to place them and the total number of permutations will be even. Hence, we can have only at most $1$ digit that appears once, and the rest must appear either 4 or 5 times ( or 0). From here we may arrive at the above $6$ configurations. Since there is an even number of odd cases, we know that the final answer will be even. After a little bit of casework we observe that 500 is far too low, and thus the answer is $\boxed{1296}.$



	\item (\textit{AMC 12 12A 2022}) A circle with integer radius $r$ is centered at $(r, r)$. Distinct line segments of length $c_i$ connect points $(0, a_i)$ to $(b_i, 0)$ for $1 \le i \le 14$ and are tangent to the circle, where $a_i$, $b_i$, and $c_i$ are all positive integers and $c_1 \le c_2 \le \cdots \le c_{14}$. What is the ratio $\frac{c_{14}}{c_1}$ for the least possible value of $r$?


 $\textbf{(A)} ~\frac{21}{5} \qquad\textbf{(B)} ~\frac{85}{13} \qquad\textbf{(C)} ~7 \qquad\textbf{(D)} ~\frac{39}{5} \qquad\textbf{(E)} ~17$

 \textbf{Answer: }E



\textbf{Solution 1}

Suppose that with a pair $(a_i,b_i)$ the circle is an excircle. Then notice that the hypotenuse must be $(r-x)+(r-y)$, so it must be the case that 
 \[a_i^2+b_i^2=(2r-a_i-b_i)^2.\] Similarly, if with a pair $(a_i,b_i)$ the circle is an incircle, the hypotenuse must be $(x-r)+(y-r)$, leading to the same equation. 

 Notice that this equation can be simplified through SFFT to 
 \[(a_i-2r)(b_i-2r)=2r^2.\] Thus, we want the smallest $r$ such that this equation has at least $14$ distinct pairs $(a_i,b_i)$ for which this holds. The obvious choice to check is $r=6$. In this case, since $2r^2=2^3\cdot 3^2$ has $12$ positive factors, we get $12$ pairs, and we get another two if the factors are $-8,-9$ or vice versa. One can check that for smaller values of $r$, we do not even get close to $14$ possible pairs.

 When $r=6$, the smallest possible $c$-value is clearly when the factors are negative. When this occurs, $a_i=4, b_i=3$ (or vice versa), so the mimimal $c$ is $5$. The largest possible $c$-value occurs when the largest of $a_i$ and $b_i$ are maximized. This occurs when the factors are $72$ and $1$, leading to $a_i=84, b_i=13$ (or vice-versa), leading to a maximal $c$ of $85$.

 Hence the answer is $\frac{85}5=\boxed{17}$.

 



\textbf{Solution 2}

Case 1: The tangent and the origin are on the opposite sides of the circle.

 In this case, $a, b > 2r$.

 We can easily prove that
 \[a + b - 2 r = c . \hspace{1cm} (1)\]

 Recall that $c = \sqrt{a^2 + b^2}$.

 Taking square of (1) and reorganizing all terms, (1) is converted as
 \[\left( a - 2 r \right) \left( b - 2 r \right) = 2 r^2 .\]

 
Case 2: The tangent and the origin are on the same sides of the circle.

 In this case, $0 < a, b < r$.

 We can easily prove that
 \[2 r - a - b  = c . \hspace{1cm} (2)\]

 Recall that $c = \sqrt{a^2 + b^2}$.

 Taking square of (2) and reorganizing all terms, (2) is converted as
 \[\left( a - 2 r \right) \left( b - 2 r \right) = 2 r^2 .\]

 Putting both cases together, for given $r$, we look for solutions of $a$ and $b$ satisfying
 \[\left( a - 2 r \right) \left( b - 2 r \right) = 2 r^2 ,\]with either $a, b > 2r$ or $0 < a, b < r$.

 Now, we need to find the smallest $r$, such that the number of feasible solutions of $(a, b)$ is at least 14.

 For equation
 \[uv = 2 r^2 ,\]we observe that the R.H.S. is a not a perfect square. Thus, the number of positive $(u, v)$ is equal to the number of positive divisors of $2 r^2$.

 Second, for each feasible positive solution $(u, v)$, its opposite $(-u, -v)$ is also a solution. However, $(u,v)$ corresponds to a feasible solution if $(a, b)$ with $a = u + 2r$ and $b = v + 2r$, but $(-u, -v)$ may not lead to a feasible solution if $(a, b)$ with $a = 2 r - u$ and $b = 2 r - v$.

 Recall that we are looking for $r$ that leads to at least 14 solutions.Therefore, the above observations imply that we must have $r$, such that $2 r^2$ has least 7 positive divisors.

 Following this guidance, we find the smallest $r$ is 6. This leads to the following solutions:

 $\left( a_1, b_1, c_1 \right) = \left( 3, 4, 5 \right)$, $\left( a_2, b_2, c_2 \right) = \left( 4, 3, 5 \right)$.

 $\left( a_3, b_3, c_3 \right) = \left( 20, 21, 29 \right)$, $\left( a_4, b_4, c_4 \right) = \left( 21, 20, 29 \right)$.

 $\left( a_5, b_5, c_5 \right) = \left( 18, 24, 30 \right)$, $\left( a_6, b_6, c_6 \right) = \left( 24, 18, 30 \right)$.

 $\left( a_7, b_7, c_7 \right) = \left( 16, 30, 34 \right)$, $\left( a_8, b_8, c_8 \right) = \left( 30, 16, 34 \right)$.

 $\left( a_9, b_9, c_9 \right) = \left( 15, 36, 39 \right)$, $\left( a_{10}, b_{10}, c_{10} \right) = \left( 36, 15, 39 \right)$.

 $\left( a_{11}, b_{11}, c_{11} \right) = \left( 14, 48, 50 \right)$, $\left( a_{12}, b_{12}, c_{12} \right) = \left( 48, 14, 50 \right)$.

 $\left( a_{13}, b_{13}, c_{13} \right) = \left( 13, 84, 85 \right)$, $\left( a_{14}, b_{14}, c_{14} \right) = \left( 84, 13, 85 \right)$.

 Therefore, $\frac{c_{14}}{c_1} = \boxed{\textbf{(E) 17}}$.

 



\textbf{Solution 3}

As $c_i$ is the length of the segment $(0,a_i)$ and $(b_i,0)$， $a_i^2+b_i^2=c_i^2$. The equation for the line that passes through $(0,a_i)$ and $(b_i,0)$ is $a_ix+b_iy-a_ib_i=0$.

 By the point-line distance formula from point $(r,r)$ to line $a_ix+b_iy-a_ib_i=0$

 \[r = \frac{ |a_ir+b_ir-a_ib_i| }{ \sqrt{a_i^2+b_i^2} }, \quad c_i = \frac{ |a_ir+b_ir-a_ib_i| }{r} = |a_i+b_i- \frac{a_ib_i}{r}|\]

 \[\because c_i \quad \text{is an integer}, \quad \therefore r|a_ib_i\]

 To find $c_1$ and $c_{14}$ for the smallest $r$, we need Pythagorean triples $(a_i,b_i,c_i)$ with relatively prime elements from the smallest. Notice that we only need to find $7$ triples with $a_i<b_i$, the $7$ other triples will simply be $(b_i,a_i,c_i)$. Using the formula $2mn, m^2 - n^2, m^2 + n^2$ to generate primitive triplets, where $m^2 + n^2$ is $c_{i},$ we have: 

 $c_{1}$ is when $m = 2, n = 1 \implies c_{1} = 2^2 + 1^2 = 5$
$c_{7} = c_{8} = ... = c_{14}$ is when $m = 7, n = 6 \implies c_{14} = 7^2 + 6^2 = 85$

 Therefore, $\frac{c_{14}}{c_1} = \frac{85}{5} = \boxed{\textbf{(E) 17}}$.

 


\end{enumerate}